%\section{Gram-determinants of one cup standard modules}
\subsubsection{Gram-determinant-decorated Rollet graph - preliminaries}
%\bajm{[needs work]} \bajm{[rename section?? this also now feels like a slightly wrong thing to emphasise in this section now]} \ppm{[agreed, so I've reorganised a tad here. No change of order, just change sectioning. Does it work?!  Another way to go would be to swap \S\ref{ss:onecup0} and \S\ref{ss:DRG1}.]}  
As noted in the Introduction (see 
\S\ref{ss:intro}),  %\ref{make-it-true!}), 
our first main aim is 
to determine the fundamental invariants of Kadar--Yu representation theory, i.e. the Cartan decomposition matrices. 
This is not easy!%
\footnote{(In the 20th century, Brauer posed the corresponding problem for the symmetric groups over fields of finite characteristic and then, coming to suspect that this problem is impossible in general, he posed the problem for the Brauer algebras over the complex field - a problem which  then remained open for over 50 years, and was only solved in the 21st century. And, viewed holistically, the Kadar-Yu problem is certainly harder.)} 
So we have some intermediate aims. 
In particular, partial data on decomposition matrices is given by morphisms from costandard to standard modules. %- or equivalently 
And partial data on these is given 
by the corresponding gram matrices 
$\Gamma(\Deltaa^{n}_{(\lambda,p)})$
of contravariant forms (in some suitable basis). 
And finally just by the determinants:
\beq
 \Delttt{n}{p,\lambda}  \; := \;  \det(\Gamma(\Deltaa^{n}_{(\lambda,p)}))\label{eq:detnot}
\eq


Even the gram determinants represent a huge swathe of complex data, raising the practical problem 
of how to organise and present what is known. 

\medskip 

%\sout{Here we collect our later results on a Rollet Graphs. Let us begin by drawing a Rollet graph for the $l=1$ case with tikz cd, in \S\ref{ss:DRG1}, ... \ppm{... or, begin with $l=-1$?...}}

\newcommand{\V}{{\mathcal V}}
\newcommand{\RollV}{\Roll^{{\mathcal V}}}
\newcommand{\RollD}{\Roll^{{\Deltaa}}}

In general a `gram-det-decorated' Rollet graph 
$\RollD_l$ 
is a Rollet graph $\Roll_l$ 
as in \ref{de:Rollet}
with the fibre of gram determinants attached to each vertex. 
And we shall include results presented in this way.
But this is quite information heavy. 
We will see from the $l=-1$ case that there the data may be simplified enormously by only attaching (for given~$l$)
the fibre of `marginal vertex functions' defined as follows. 
The marginal vertex function is 
%\bajm{[makes sense to me to make this a function of $m$ not $n$]} \ppm{[I am wary of making a change late in the day (and I don't see why yet). - makes yet another variable for the reader to remember. Every other such thing uses the same set of superscript and subscript! To discuss.]}\bajm{[okay]}
%
%\newcommand{\V}{{\mathcal V}}
\beq  \label{eq:MVF1}
\V^{\bajmm{n}}_{(p.\lambda)} \; \bajmm{:}= \; 
\frac{ \Delttt{(n=p+2m)}{p,\l} }{ \prod_{{v\; adjacent\; \atop to \; (p,\lambda)\; in\;\Roll_l }} \Delttt{(n-1)}{v}}
\eq 
- recall that $p$ is the number of propagating lines.  
%\ppm{[in (\ref{def:indexset}) we write $(\lambda,p)$. which is it?!]}\bajm{[i propose $(p,\l)$ it used more throughout this section and the original KY ref]}\ppm{[yes.]}
In the new decorated graph the fibre of determinants is replaced by this $\V$ as $m$ varies - if indeed it depends on $m$.
%\[ ... \]
\ppmm{So we will also include some results presented in this way.} 
Let us write $\RollV_l$ for this particular decorated Rollet graph. 

\mdef   \label{de:Roll-etc}
Finally, and even less information-heavy, we may also present a certain section through $\RollV_l$,
corresponding to fixed $m$ 
(fixed number of `cups' in the diagrams of the diagram basis) - which we denote by $\Roll^{(m)}_l$.

The corresponding section through $\RollD_l$ is denoted $\Roll^{m}_l$. 

\medskip 

Observe that if we can determine the marginal vertex functions then we can determine the \bajmm{G}ram determinants by a recursion. %\bajm{[in my reading Gram is with capital G as it is a name?]}\ppm{[probably right. can conform if time. (if not, there is always fourier!)]}
So a next step is to study the marginal functions. 
In order to give them in closed form it will be convenient to use Chebyshev polynomials. 
And in particular this brings us to the `bootstrap' one-cup cases. 
So in 
\S\ref{ss:onecup0} and in particular in 
\S\ref{ss:cheby0} we collect the manipulations that we will need. 

Then in \S\ref{ss:e-next!} we will address the general case. 
%.... \bajm{[finish this?]}\ppm{[or just kill ellipsis for now.]}



\section{Gram-determinants of one cup standard modules}  \label{s:onecup0}

In this section, we consider the gram determinants of the $J_{l,n}$ standard modules $\Deltaa_{(n-2,\l)}^{n}$ for $n\geq l+4$. The main result is Theorem \ref{thm:main} in \S\ref{ss:non1dsp} which allows us to obtain closed forms for $\Deltaa_{(n-2,\l)}^{n}$ as a series in $n$ by computing the first two cases (\textit{i.e.} $n=l+4,l+5$); these cases uniquely determine a Chebyshev series of polynomials, $P^{(\l)}_{n}$, with which we can express each determinant. We apply this theorem to compute closed forms of this determinant for all $\l\vdash l+2$ with $l\leq 2$ and additionally in the cases $\l=(l+1,1), (l+1,1)^T$ for $l\leq 7$. The polynomials $P^{(\l)}_{n}$ exhibit some remarkable properties which we begin to document in \S\ref{ss:examples0}. We study the distribution of roots of these polynomials in \S\ref{s:roots}.

\subsection{On Chebyshev polynomials (for gram determinants)}  \label{ss:cheby0}

%Some preparation: 
%\\
%\sout{\ppm{[I don't really understand this formulation yet:]}
%Recall, for a series of functions, $P=\{ P_n\}_{n\in \Z}$, the Chebyshev recurrence is} 
%\\
\bajmm{We say that a series of functions, $P=\{ P_n\}_{n\in \Z}$, is a ``Chebyshev series" if it satisfies the Chebyshev recursion:}
\beq \label{eq:Cheby0} 
P_{n+1}(x) = x P_{n}(x) +P_{n-1} (x)   ,
\eq  
(see e.g. \cite[\S6.3.3]{Martin91} and references therein; cf. also e.g. \cite{KarlinStudden}).
%\ppm{\sout{When a series $P$ obeys the recurrence it is a Chebyshev series.}}
This means of course that the series is 
determined in principle by any two successive values, $P_i$ and $P_{i+1}$. 
Classically these would be with $i=0$, as we will recall shortly. 
Two series $P,P'$ are `parallel' if $P_n = P'_{n+i}$ for some $i$. 

%\ppm{[some examples here? could have classical and one where $i$ below is not 1, say?]}

\mdef  \label{pr:cheb1}
Observe that if any two successive functions $P_i$ and $P_{i+1}$ are (integral) polynomial then all $P_n$ are. If in addition $\deg(P_{i+1})=\deg(P_i) + 1$ \ppmm{for some $i$,} then $\deg(P_{k+1})=\deg(P_{k}) + 1$ for all $k>i$ 
%\ppm{[-this appears to have different nature than first one, since only claimed to ascend.]}\bajm{[fixed]}
%\\
and further, if in addition
%\ppm{[-meaning the degree is ascending?...]}\bajm{[yes]}
$P_{i+1}$ is monic 
%\ppm{[...in which case is it enough that $P_{i+1}$ is monic?]}\bajm{[true]}
then so is $P_k$ for all $k>i$. 

\mdef We call a Chebyshev series of polynomials $P$ ``reduced" if there is no common non-trivial (\textit{i.e.} degree $>0$) polynomial factor 
%\sout{\ppm{[what ring?]}}
among all $P_i$ (equivalently, among any two successive %\sout{\ppm{non-zero [????]}\bajm{[no actually, we don't want non-zero here]}} \ppm{[I don't understand. if a term is zero then the whole next term is common, but only for those two.]}\bajm{[not quite, that term (the first non-zero one) will be a factor of the whole series (if $P_n=0$ then the cheb. rec. implies that $P_{n+2}(x)=x P_{n+1}(x)$) so the statement as is holds - i.e. the only reduced series with a 0 term somewhere is the classical one]}
terms). 
\\ 
Any polynomial series $P$ is a (polynomial) multiple %\ppm{[-is proportionality normally by a constant?]}
of a reduced one $P'$; $P'$ can be obtained from $P$ by removing a maximal common factor among two successive terms. 

%\ppm{[maybe reorganise a tad here. just for readability. R2 etc]}

\mdef  \label{de:Chebyab}
Let us write $\PP^{a,b} $ for the Chebyshev series with initial conditions 
$\PP^{a,b}_0 = a$ and $\PP^{a,b}_1 = b$.  

We may abuse notation slightly and, provided $i$ is fixed, write $\PP^{a,b} $ 
for the parallel series that starts with $\PP^{a,b}_i =a$. 

%\sout{Observe that if any two successive terms are polynomial then all are. 
%\\
%And if $\deg(b)=\deg(a)+1$ then $\deg(\PP^{a,b}_{k+1} = \deg(\PP^{a,b}_k)+1$ for all $k>i$. }
%\ppm{[its getting quite long. could delete some of the chatty notes-to-self? (nice though they are):]}
As we will see, explicit usage of the $\PP^{a,b}$ notation is often cumbersome in practice for general series (since $a$ may be a lengthy expression), and we will largely restrict the notation to abstract manipulations. Instead we use more specific notations for interesting cases, as in \S\ref{ss:onecup0}.  


\mdef 
The classical initial conditions are $P_0(x) =0$ and $P_1(x) =1$ giving then 
$P_2(x) =x $; $P_3(x) = x^2-1$; $P_4(x) = x(x^2 -2)$; and so on. In our $\PP^{a,b}$ notation this classical series is $\PP^{0,1}$. 

\mdef 
For convenience we set $P^U:=\PP^{0,1}$ since this series is related to the classical series of ``U-type" Chebyshev polynomials, $U$, which can be defined by the identity
\beq U_n(\cos(\th))\sin(\th)= \sin((n+1)\th),\label{eq:trig}\eq
by $P^U_n(x)=U_{n-1}(x/2)$. The polynomials $U_n$ and $P^U_n$ are well defined for $n<0$ with the above formulae still holding here; indeed, one can use these to determine
\beq P^U_{-n}(x)=-P^U_{n}(x),\label{eq:negcheb} \eq
for any $n\in \Z$. 

\mdef   \label{pr:Cbasis} Note that $\{ P^U_1(x),P^U_2(x),\dots \}$ is a $\C$-basis for the polynomial ring $\C[x]$.

\ignore{{ %%% all these promoted now to macros/macros-main.tex   %%%
%\newcommand{\ptarmigan}{Cheby condition}
\newcommand{\ptarmigan}{Cheby}
\newcommand{\ramping}{ramping}
\newcommand{\Ramping}{Ramping}
}} 
%I didn't see this command, I have replaced "Chebyshev condition" everywhere with "\bajm{X} property". Once we decide an appropriate name we can just globally replace \bajm{X} everywhere... this won't cause any issues
%LOL


\subsubsection{The \Ramping\ Property}

%\ppm{[-Ha, I don't think I understand killing the macro though. why would one replace a macro that can instantly and uniformly be changed to write anything you decide, with a passive symbol that will have to be edited across multiple instances!?]}\bajm{[you are correct I didn't see it, still a global replacement of ``bajm{X}" should do the job]}

Here we discuss a property that Chebyshev series may 
achieve eventually. 

%not quite sure how to proceed now. I guess I'll try again. 
\newcommand{\rampsfrom}{ramps from}   %%a ramp being ascending at, plausibly, 45 degrees.  :-) 
% OK you don't like this either. never mind. I've added the \ramping macro above.


\mdef\label{def:chebycond} 
Let $P$ be a 
%\sout{\ppm{reduced (hence polynomial) [????]}}\bajm{[I think it is better to let these be part of the definition not preconditions so I get more mileage by calling something X...] } 
Chebyshev series. For some integer $N$, we say that the series $P$ 
%\ppm{\em{\rampsfrom} $N$ [????]}\bajm{[decent name, but it would be much easier if we can just find a replacement for X since the rest of the paper can then be fixed by a global replacement]}
has the ``\ramping\ property 
%\ppm{[-any reason for this name?]}\bajm{[No (other than that the classical chebyshev series give an example) I just called it something to avoid wasting time choosing a name (which I have a tendency to do), happy for a better name :)]}
for $n\geq N$'' if it is polynomial and reduced, and if, furthermore, for $n\geq N$, $P_n$ is monic and $\deg(P_{n+1})=\deg(P_{n})+1$. 
%\ppm{[is reduced not an all n condition? and polynomial?]}\bajm{[it is, I am trying to say the last two conditions only apply for $n\geq N$]}
%\ppm{[as defined I think \ramping\ for N implies ramping for N+1 and so on. do you want to make N the minimum one? we could say:]  Series $P$  {\em \rampsfrom} $N$ if $N$ is minimal. [?]}\bajm{[No i actually don't because the polynomials we see later (the $P_n^{(\l)}$) appear to ramp for $n\geq l+3$ but I can't see a way to prove this generally - whereas we have a proof for $n\geq l+4$]}
Such a series is uniquely determined by $N$ and the ratio $P_{N+1}(x)/P_{N}(x)$, a rational function.




%\sout{Here the series is determined not just in principle, but a simple re-parametrisation gives it in a useful closed form, that we will recall shortly.  }
%\sout{In our $\PP^{a,b}$ notation this classical series is $\PP^{0,1}$. }

\mdef \label{pa:chebslv} 
Suppose that a series of polynomials $P$ has the \ramping\ property for $n\geq N$, and suppose that $d=\deg(P_N)$. 
%\sout{We claim  \ppm{[? - not sure what's going on here]} that one can determine} 
Then there exist unique coefficients $a_k, b_i, b'_j \in \C$, 
%for $k=-(d+1),-d,\dots,d,d+1$, $i=1,\dots, d+1$, and $j=1,\dots,d+2$
such that
\begin{align*}
    P_N(x)&=\sum_{k=-(d+1)}^{d+1} a_{k} P^U_{k}(x)=\sum_{i=1}^{d+1} b_i P^U_{i}(x),\\
    P_{N+1}(x)&=\sum_{k=-(d+1)}^{d+1} a_{k} P^U_{k+1}(x)=\sum_{j=1}^{d+2} b'_j P^U_{j}(x).\end{align*}
%\ppm{[-what is $b_i$? hark back to \ref{pr:Cbasis} maybe?]}\bajm{[okay?]}
The $a_k$ can be determined as follows: first solve for the coefficients $b_i$ and $b_j'$ - these are unique by \ref{pr:Cbasis}. Then by using \eqref{eq:negcheb} we obtain the following full rank linear system relating the $a_k$ to the $b_i$ and $b_j'$:
\begin{align*}
    b_k&=a_k-a_{-k}, & b'_{d+2}&=a_{d+1}, & b'_{d+1}&=a_{d}, &b'_{k'}&=a_{k'-1}-a_{-(k'+1)},
\end{align*}
for $k=1,\dots,d+1$ and $k'=1,\dots, d$. In particular, monicity of $P_N$ and $P_{N+1}$ imply that $a_{d+1}=1$ and $a_{-(d+1)}=0$.

Since the relation \eqref{eq:Cheby0} is linear, it thus follows that for any $l\in \Z$, we have
\beq 
P_{N+l}(x)=\sum_{k=-(d+1)}^{d+1} a_{k} P^U_{k+l}(x)=P^U_{d+1+l}(x)+\sum_{k=-d}^{d} a_{k} P^U_{k+l}(x). 
\eq




\ignore{
\mdef 
If we start with $P'_2(x)=2x$ and $P'_3(x) = (x-1)(x+2) \;$ 
(see e.g. \cite[\S5.2]{ShBr})
then  $P'_4(x) = x(x^2+x-4)$,
\beq   \label{eq:l=0+}
P'_5(x) = x^4 + x^3 -5x^2 -x +2, \qquad 
P'_6(x) =  x(x-1)(x^3 +2x^2 -4x -6)   \qquad 
\eq  
and so on. 
A better notation is perhaps $P'_n(x) = P^{(2)}_n(x)$, 
since this sequence is related to the $\lambda = (2)$ branch - as we will explain shortly. 
\\ \ppm{
Suppose now we start with $P''_2(x) = 2$ and $P''_3(x) = x-1$  
- a better notation here is $P''_n(x)  = P^{(1^2)}_n(x)$.
Then ... }

\mdef  \label{pa:Cheby1}
Here are tables of values for various initial conditions 
(we say `initial' conditions, but such Chebyshev sequences are determined by any two successive terms; other than for $P$, the starting point for $n$ values is not yet locked down):
\[
\begin{array}{c|c|c|ccc}
n\setminus  & P_n(x)        & P'_n =P^{(2)}_n      & P''_n = P^{(1^2)}_n \\
2      &  x                 & 2x               & 2 \\
3      & x^2 -1             & (x-1)(x+2)        & x-1  \\
4      & x(x^2-2)           &  x(x^2+x-4)      & (x+1)(x-2)  \\ 
5      & (x^2-x-1)(x^2+x-1) & x^4+x^3-5x^2 -x +2  & (x^3-x^2-3x+1)  \\
6       & x(x^2-1)(x^2-3)   & x(x-1)(x^3+2x^2-4x-6) & (x-1)(x^3-4x-2) 
\end{array}
\]
and continuing (with some initial conditions suggested by $l=1$, cf. \ref{pa:l1}):
\[ \hspace{-.861cm} 
\begin{array}{c|c|c|ccc}
n\setminus  & P^{(3)}_n             & P^{(21)}_n      &  P^{(1^3)}_n \\
1    &    &  3(x^2-1) & \\
2      &  3(x+2)             & x(x^2-4)               & 3 \\
3      &    x(x+4)            & (x^4-7x^2+3)         & x-2  \\
4      & (x+1)(x^2+3x-6)        &  x(x^2-1)(x^2-7)      & (x+1)(x-3)  \\ 
5      & x(x^3+4x^2-4x-10)        & (x^6-9x^4+14x^2-3)  & (x^3-2x^2-4x+2)  \\
6       & (x^5+4x^4-5x^3-14x^2+3x+6) & %\ppm{??(x^7-...x^6+...)} 
x \left(x^{6}-10 x^{4}+22 x^{2}-10\right)
& (x^4-2x^3-5x^2+4x+3) \\ 
7   &  \dots  & 
%\left(x -1\right) \left(x +1\right) \left(x^{6}-10 x^{4}+21 x^{2}-3\right)
\left(x^2 -1\right) \left(x^{6}-10 x^{4}+21 x^{2}-3\right)
& \dots 
\end{array}
\]
%\ppm{[-so far  unfinished! ]}



%\newcommand{\PP}{{\mathtt P}}  %% hook for Chebys 


\mdef 
In general there is a Chebyshev sequence for each `branch' for each $l$, 
hence for each $\lambda\vdash l+2$, 
thus a sequence $\PP^\lambda$ 
\ppm{defined specifically by ...
}

The first perspective on this is that the gram matrices on a branch are `related' by the Chebyshev recursion. This is true on the nose if we look at the slice through the Bratelli diagram corresponding to 1-cup modules $\Deltaaa{n}{(n-2,\lambda)}$ with $\lambda\in \{ (3), (1^3) \}$. 
... }






\subsection{Gram Determinants of One Cup Standard Modules}\label{ss:onecup0}
Here we consider the Gram determinants of the standard modules $\Deltaa_{(p, \l)}^{n}$ for $n=p+2$. 

\subsubsection{The one-cup gram matrix}    \label{ss:gram-mat}

%\ppm{In this section we ...}

As in \cite[Prop. 4.22]{KadarYu}, we may construct a basis for 
the left module 
$\Deltaaa{n}{(p, \l)}$ 
for given $l$ 
from a basis for the corresponding Specht module, $\mathcal{S}^{\l}$, and the set of half diagrams $J^{||}_{l}(n,p)$ (\textit{cf.} \ref{de:stdmods}). 
Indeed this holds for any $p$, but if $p=n-2$ as here then the latter half diagrams are completely determined by the position of a single ``cup" hence our referring to the $\Deltaa_{(p, \l)}^{n}$ as ``one cup modules". 

\mdef   \label{de:uu} 
For fixed $l$ and $n$, write $\uu_{ij}$, with $i<j$ for the half diagram $\uu_{ij}\in J_l^{||}(n,n-2)$ 
with a cup joining the sites $i$ and $j$.  
%\bajm{[is the $e_{ij}$ notation okay?]} \ppm{[yes. make it a macro just in case we want to change the font away from default? (I might have tried ${\mathsf u}_{ij}$ just because $u$ looks more like a cup and less like an identity, but no biggie.)]}\bajm{I will switch to ${\mathsf u}_{ij}$ because I think its better!} 
%have made a macro \uu
For example, when $n=7$ we have 
\[
\uu_{24} =  \onecupp{8}{2}{4}{1}{3}{5}{6}{7}      \hspace{.31cm} \in J_{1}^{||}(7,5). 
\]
Note that if $j=i+1$ then $ht(\uu_{ij})=-1$, otherwise $ht(\uu_{ij})=j-3$.

%\ppm{[say something about height $l$? or drop the $l$ here?]}\bajm{[the height is implicit in the notation $J_l$...?]}\ppm{[my point exactly.]}\bajm{[ah gotcha]}

\newcommand{\BBB}{{\mathsf B}}  %% hook for diagram bases

\mdef 
Write $\BBB_{(l,n)}$ as shorthand for the set of half diagrams $J_{l}(n,n-2)$. One has that 
\beq
\BBB_{(l,n)}=\{\uu_{12},\uu_{23},\dots,\uu_{n-1,n}\}\cup \left(\bigcup_{k=3}^{\min{n,l+3}} \BBB_k\right), \label{cupbasis}
\eq
where $\BBB_k:=\{\uu_{j,k} \ | \ j=1,\dots, k-2\}$. 
Observe that $\BBB_k$ is precisely those one-cup half-diagrams with height $k-3$. 

\mdef  \label{pa:cuprec}
Observe for $n\geq l+3$, we have that $\BBB_{(l,n+1)}= \BBB_{(l,n)}\cup \{\uu_{n,n+1}\}$.

\mdef 
Fix $l$ and $\l \vdash l+2$. When $n\geq l+4$, consider a basis $\mathcal{B}=\{b_1,\dots, b_{d_{\l}}\}$ for 
$\mathcal{S}_{\l}$ regarded as one for $\Delta_{(n,\l)}^{n}$ (\textit{c.f.} \ref{pa:spbasis})
%\ppm{[-there is no action of J on S, so better to write as Delta or keep saying iso is as sym-module (which maybe seems a bit odd) or as vector space.]}
with $b_i=x_i c_{\l}$ for some $x_i\in \C \Sym_{l+2}\subset J_{l,n}$ and $c_{\l}$ as per \eqref{de:cl}. Then a basis for $\Deltaaa{n}{(n-2,\l)}{}$ can be obtained as 
\beq
\BBB^n_{(n-2,\l)} \; := \; \{ \uu_{jk}b_i \ | \ \uu_{jk} \in \BBB_{(l,n)}, \ b_i \in \mathcal{B}  \}. \label{eq:basis} \eq

\mdef \label{pa:gramrels}
For $n\geq l+4$, we now consider the form $\langle\_,\_\rangle$ on $\Deltaa^{n}_{(n-2,\l)}$. The value of the form $\langle \uu_{jk} b_i , \uu_{j'k'} b_{i'} \rangle$ is determined by the equation
 \[
 c_{\l }   x_{i'}^* \uu_{j'k'}^* \uu_{jk} x_i  c_{\l}= \langle b_i \uu_{jk}, b_{i'} \uu_{j'k'} \rangle c_{\l}.
\]
There are three cases to consider, namely $|\{j,k,j',k'\}|=2,3,4$. In the first case $ \uu_{j'k'}^*\uu_{jk}=\a e$, in the second $\uu_{j'k'}^* \uu_{jk} =\sigma$ for some $\sigma\in \Sym_{l+2}$. In the last case,  $\uu_{j'k'}^* \uu_{jk} $ has $p-2$ propagating lines. We thus find
\[ \langle  \uu_{jk} b_i ,\uu_{j'k'} b_{i'}  \rangle =\begin{cases}
    \a \langle x_i, x_{i'}\rangle, & |\{i,j,i',j'\}|=2,\\
    \langle \sigma x_i, x_{i'}\rangle, & |\{i,j,i',j'\}|=3, \text{ where } \sigma=\uu_{j'k'}^*\uu_{jk} ,\\
    0, & \text{otherwise.}
\end{cases} 
\]
where $\langle\_,\_\rangle$ also here denotes the form on $\Delta^{n}_{(n-2,\l)} c_{\l} \simeq \SS_\l$ (\textit{c.f} \ref{de:stdmods}). Passing to an orthonormal basis for $\SS_\l$, with respect to this form, we observe that $\a$ only appears on the leading diagonal of the matrix. In particular, $\Gamma(\Delta^{n}_{(n-2,\l)})$ now has the form
\beq \Gamma(\Delta^{n}_{(n-2,\l)})(\a)=\a \, \text{I} +\Gamma(\Delta^{n}_{(n-2,\l)})', \label{eq:adep} \eq
where $\Gamma(\Delta^{n}_{(n-2,\l)})'=\Gamma(\Delta^{n}_{(n-2,\l)})|_{\a=0}$ is a real symmetric matrix. As observed for example in \cite{KY2}, this implies that the form $\langle \_ , \_\rangle$ is non-degenerate for $\a$ away from the real line.

For $l+3<j<n$ basis elements of the form $\uu_{j,j+1} b_i $ 
%ppm{[-wrong way round? find and fix all such cases?]}\bajm{[good catch]}
only have non-zero pairing with those of the form $\uu_{j',j'+1} b_{i'}$ for $j'=j,j\pm 1$, where we may compute
\beq \langle \uu_{j,j+1} b_i  , \uu_{j',j'+1}b_{i'} \rangle = \delta_{ii'}(\a \delta_{jj'}+\delta_{j+1,j'}+\delta_{j-1,j'}). \label{cheby} \eq

Equation \eqref{cheby} is the source of an important property of our Gram matrices $\Gamma(\Deltaa^{n}_{(p,\lambda)})$. 

\subsubsection{When $\SS_\l$ is one dimensional}\label{ss:1dspecht}

The underlying theory for this case is treated in \cite{KadarYu,ShBr}. 
This section provides a review in preparation for the general case in \S\ref{ss:non1dsp}, 
and examples extended beyond those previously computed. 

\medskip 

Suppose firstly 
then, 
that $\SS_\l$ is one dimensional, \textit{i.e.} $\l=(p)$ or $(1^p)$, and so $\Deltaa_{(p,\l)}^{n}$ has basis 
$\{ \uu_{ij} c_{\l} \}_{ij}$ with $\uu_{ij}\in B_{(n,l)}$. 
%\ppm{[maybe there is just a bit too much talking, about the 1d case? that is already treated in the literature, certainly more than the higher d. we only want as much of it as helps to set the scene? not sure. too much is better than too little.]}
Furthermore, it is naturally divided into two ``regimes" whereon the form $\langle \_ , \_ \rangle$ displays distinct behaviour; the first we term the ``head" consisting of basis elements $ \uu_{ij} c_{\l}$ with $j\leq l+4$, and the second the ``tail" consisting of those same elements with $j>l+4$. The gram matrix restricted to the head is a dense matrix, whereas for the tail it is tridiagonal (when appropriately ordered). Furthermore, the pairing between elements of the head and the tail only has one non-zero case, namely $\langle \uu_{l+3,l+4} c_{\l}, \uu_{l+4,l+5} c_{\l} \rangle=1$.

\mdef 
To illustrate this behaviour, we give the gram matrix for the $J_{1,13}$-module $\Deltaa^{(13)}_{(11,\lambda)}$, for $\l =(3)$ or $\l=(1^3)$ 
%below 
(by taking $\e=1$ or $\e=-1$, respectively)
below:
%\ppm{The indicated partitioning is into head and tail:}
\beq \left(
\begin{array}{ccc ccc c| cccccccc}
 \alpha  & 1 & \e & 1 & \e & 0 & 0 & 0 & 0 & 0 & 0 & 0 & 0 & 0 & 0 \\
 1 & \alpha  & 1 & 1 & 0 & \e & 0 & 0 & 0 & 0 & 0 & 0 & 0 & 0 & 0 \\
 \e & 1 & \alpha  & 0 & 1 & \e & 1 & 0 & 0 & 0 & 0 & 0 & 0 & 0 & 0 \\
 1 & 1 & 0 & \alpha  & 1 & 1 & 0 & 0 & 0 & 0 & 0 & 0 & 0 & 0 & 0 \\
 \e & 0 & 1 & 1 & \alpha  & 1 & \e & 0 & 0 & 0 & 0 & 0 & 0 & 0 & 0 \\
 0 & \e & \e & 1 & 1 & \alpha  & 1 & 0 & 0 & 0 & 0 & 0 & 0 & 0 & 0 \\
 0 & 0 & 1 & 0 & \e & 1 & \alpha  & 1 & 0 & 0 & 0 & 0 & 0 & 0 & 0 \\\hline
 0 & 0 & 0 & 0 & 0 & 0 & 1 & \alpha  & 1 & 0 & 0 & 0 & 0 & 0 & 0 \\
 0 & 0 & 0 & 0 & 0 & 0 & 0 & 1 & \alpha  & 1 & 0 & 0 & 0 & 0 & 0 \\
 0 & 0 & 0 & 0 & 0 & 0 & 0 & 0 & 1 & \alpha  & 1 & 0 & 0 & 0 & 0 \\
 0 & 0 & 0 & 0 & 0 & 0 & 0 & 0 & 0 & 1 & \alpha  & 1 & 0 & 0 & 0 \\
 0 & 0 & 0 & 0 & 0 & 0 & 0 & 0 & 0 & 0 & 1 & \alpha  & 1 & 0 & 0 \\
 0 & 0 & 0 & 0 & 0 & 0 & 0 & 0 & 0 & 0 & 0 & 1 & \alpha  & 1 & 0 \\
 0 & 0 & 0 & 0 & 0 & 0 & 0 & 0 & 0 & 0 & 0 & 0 & 1 & \alpha  & 1 \\
 0 & 0 & 0 & 0 & 0 & 0 & 0 & 0 & 0 & 0 & 0 & 0 & 0 & 1 & \alpha  \\
\end{array}
\right).\eq
The basis $\uu_{ij} c_{\l}$ is ordered lexicographically, 
thus
\[
\hspace{-.5cm}
\onecupp{8}{1}{2}{3}{4}{5}{6}{7} ,
\onecupp{8}{1}{3}{2}{4}{5}{6}{7} ,  
\onecupp{8}{1}{4}{2}{3}{5}{6}{7} ,  
\onecupp{8}{2}{3}{1}{4}{5}{6}{7} , \cdots
\]
The lines in the matrix are separating the head and the tail.  

\mdef 
More generally, for any $n>l+4$, and $\l=(l+2), (1^{l+2})$, by ordering the basis in the same way the matrix $\Gamma(\Deltaa^{n+1}_{(n-1,\lambda)})$ has the form
\beq
\begin{pNiceArray}{cccc|cc}
\Block{4-4}<\large>{\Gamma(\Deltaa^{n-1}_{(n-3,\lambda)})} &&&& 0&0\\
&&&& \vdots & \vdots\\
&&&& 0 & 0 \\
&&&& 1 &0\\ \hline
0 & \dots &0 &1 & \a &1 \\
0 & \dots &0 &0 & 1 &\a \\
\end{pNiceArray}.
\eq

%\subsubsection{The gram determinant}  % when the Specht module is 1d}

To compute the determinant $\Delttt{(n+1)}{n-1,\lambda}$ for such $\l$ then,
we can expand along the last row obtaining the recursive formula
\begin{align}\Delttt{(n+1)}{n-1,\lambda}(\a)&= \a  \ \left| \ \begin{NiceArray}{cccc|c}
\Block{4-4}<\large>{\Gamma(\Deltaa^{(n-1)}_{(n-3,\lambda)})} &&&& 0\\
&&&& \vdots \\
&&&& 0  \\
&&&& 1 \\ \hline
0 & \dots &0 &1 & \a 
\end{NiceArray}\right| \ - \  \left| \ \begin{NiceArray}{cccc|c}
\Block{4-4}<\large>{\Gamma(\Deltaa^{(n-1)}_{(n-3,\lambda)})} &&&& 0\\
&&&& \vdots \\
&&&& 0  \\
&&&& 0 \\ \hline
0 & \dots &0 &1  &1 \\
\end{NiceArray}\right | \nn \\ 
&=\a\, \Delttt{(n)}{n-2,\lambda}(\a)-\Delttt{(n-1)}{n-3,\lambda}(\a).\label{eq:detrec1}
\end{align}
That is, the determinants, $\Delttt{(n)}{n-2,\lambda}$, satisfy the Chebyshev recursion in $n$ \eqref{eq:Cheby0}. In other words, from \S\ref{ss:cheby0} we have
\[
\Delttt{(n)}{n-2,\lambda}   = \PP^{a,b}_n
\]
with initial conditions $a=\Delttt{(l+4)}{l+2,\lambda}$ and $b=\Delttt{(l+5)}{l+3,\lambda}$ at $n=l+4$ and $n=l+5$ respectively. 
%\ppm{[-using some parallel-ness?]}\bajm{[sure - can you think of how to word this?]}
We will apply this in \S\ref{ss:examples0} to give explicit formulas for all Gram determinants in the cases $\l=(l+2), (1^{l+2})$ for $0\leq l \leq 12$.

\ignore{\subsubsection{\ppm{(temporary sec. marker)} Applications and examples}
\bajm{[I agree this may be more natural]}
\ppm{[could move this subsec. to after Theorem to avoid inverted referencing. not sure.]} \ppm{[-could even integrate into the main such section now?]}

\mdef \label{pa:1dgdets}
For $l=0$  %, we consider the row and column partitions $\l$ here.  In these cases, 
direct computation yields
\begin{align*}
\Delttt{(4)}{2,(2)}&=(\alpha -1) \alpha  \left(\alpha ^2+\alpha -4\right),& \Delttt{(5)}{3,(2)}&=(\alpha -1) \left(\alpha ^4+\alpha ^3-5 \alpha ^2-\alpha +2\right),\\
    \Delttt{(4)}{2,(1^2)}&=(\alpha -2) (\alpha -1) (\alpha +1) (\alpha +2),&
   \Delttt{(5)}{3,(1^2)}&=(\alpha -1) (\alpha +2) \left(\alpha ^3-\alpha ^2-3 \alpha +1\right),
\end{align*}  
For $l=1$, we consider the row and column partitions $\l$ here. 
In these cases, direct computation yields
\begin{align*}
%   \\
    \Delttt{(5)}{3,(3)}&=(\alpha -2)^2 \alpha ^2 (\alpha +1) \left(\alpha ^2+3 \alpha -6\right),& \Delttt{(6)}{4,(3)}&=(\alpha -2)^2 \alpha ^3 \left(\alpha ^3+4 \alpha ^2-4 \alpha -10\right),\\
    \Delttt{(5)}{3,(1^3)}&=(\alpha -3) (\alpha -2)^2 (\alpha +1) (\alpha +2)^3,&
    \Delttt{(6)}{4,(1^3)}&=(\alpha -2)^2 (\alpha +2)^3 \left(\alpha ^3-2 \alpha ^2-4 \alpha +2\right).
\end{align*}
Thus, using the notation of Theorem \ref{thm:main}, 
\ppm{[-at the mo this is a forward ref!]}
we have here $\Delttt{(n)}{n-2,\l}=C^{(\l)}(\a) P^{(\l)}_{n}(\a)$ where
\begin{align*}
    C^{(2)}(\a)&=(\a-1), &    C^{(1^2)}(\a)&=(\a-1)(\a+2),\\
    C^{(3)}(\a)&=(\a-2)^2\a^2, &    C^{(1^3)}(\a)&=(\a-2)^2(\a+2)^3,
\end{align*}
and the sequences $P_n^{(\l)}(\a)$ are completely determined by the Chebyshev recursion, and the ``initial conditions":
\begin{align*}
    P_4^{(2)}(\a)&=\a(\a^2+\a-4), &        P_5^{(2)}(\a)&=\a^4+\a^3-5\a^2-\a+2,\\
    P_4^{(1^2)}(\a)&=(\a-2)(\a+1), &        P_5^{(1^2)}(\a)&=\alpha ^3-\alpha ^2-3 \alpha +1,\\
    P_5^{(3)}(\a)&=(\alpha +1) \left(\alpha ^2+3 \alpha -6\right), &        P_6^{(3)}(\a)&=\alpha  \left(\alpha ^3+4 \alpha ^2-4 \alpha -10\right),\\
    P_5^{(1^3)}(\a)&=(\alpha +1) \left(\a-3\right), &        P_6^{(1^3)}(\a)&=  \left(\alpha ^3-2 \alpha ^2-4 \alpha +2\right).
\end{align*}

\mdef   \label{pa:able}
By making use of \ref{pa:chebslv} we may write each polynomial $P_n^{(\l)}$ explicitly. 
We do this in more generality; we are able to compute 
\ppm{[-let's reword to show the proof.]}\bajm{[what this means?]}\ppm{[ah I see, more generality does not mean generality. ok I'll just think of a wording a bit different from ``are able to''.]}
(also in the notation of Theorem \ref{thm:main}) 
\ppm{[-forward ref.]}
for $0\leq l\leq 12$, that
\begin{align*}
    C^{(l+2)}(\a)&=(\a+(l-1))^{l+1}(\a-2)^{\frac{1}{2}l(l+3)},\\
    C^{(1^{l+2})}(\a)&=(\a-(l+1))^{l+1}(\a+2)^{\frac{1}{2}(l+2)(l+1)},
\end{align*}
and
\begin{align*}P^{(l+2)}_{l+4+k}&=P_{k+4}^U(\a)+(3 l+1) P_{k+3}^U(\a)+
\left(2 l^2-l-2\right)P_{k+2}^U(\a)\\
&\quad+(l+1)((2l-1)P^U_{k+1}(\a)-P^U_{k}(\a)),\\
P^{(1^{l+2})}_{l+4+k}&=P_{k+3}^U(\a)-(l+1)(P_{k+2}^U(\a)+P_{k+1}^U(\a)).\end{align*}
Observe in each case:
\begin{align}
    P^{(l+2)}_{l+4}&=(\alpha +l) \left(\alpha ^2+(2l+1)\alpha -2(l+1)\right), \label{eq:row0}\\
    P^{(1^{l+2})}_{l+4}&=(\alpha +1) (\alpha- (l+2) ). \label{eq:col0}
\end{align}}



\subsubsection{On Gram determinants $\Delttt{(n)}{n-2,\l}$ 
for general $\l$
%when $\SS_\l$ is not one dimensional
}\label{ss:non1dsp}

%\mdef  
When $\SS_\l$ is not (necessarily) one dimensional we may proceed as follows. 
In this section we give a formal factorisation of the gram determinants. 
In \S\ref{ss:examples0} we determine the factors. 


\mdef 
Assume that $\BB=\{b_1,\dots, b_{d_\l}\}$ is an orthonormal basis for 
$\SS_\l$, regarded as one for $\Delta^n_{(n,\l)}$ (\textit{cf.} \ref{pa:spbasis}),
%\ppm{[-same issue yet again.]}
with respect to our contravariant form $\langle \_ , \_ \rangle$. Observe that this requires that we are working over a splitting field for the symmetric group. And in fact we will go further and work here over $\C$. More precisely we are working over $\C[\xx]$ so we cannot, for example, take an orthonormal basis for $\Deltaa^{n}_{(p,\l)}$.

The formula \eqref{cheby} indicates that we can divide our basis for $\Deltaa^{n}_{(p,\l)}$ into a head and a tail for each basis element $b_i \in \BB$, with similar behaviour. 
In this case, there may be some non-zero values of the form between elements belonging to distinct heads. 
For example, when $\l=(2,1)$, by choosing the basis 
\beq   \label{eq:21ortho}
{\mathcal B}_{(21)} \; 
   = \;\;\{b_1=c_{\l}, \; b_2=\frac{2}{\sqrt{3}}((23)+\frac{1}{2} e )c_{\l}\}
\eq  
for $\SS_\l$, and ordering the 
$ \uu_{j,k} b_i$ 
in lexicographic order now with respect to the triple $(i,j,k)$, we obtain the %following 
expression for the gram matrix $\Gamma(\Deltaa^{8}_{(6,\l)})$ 
in Fig.\ref{fig:n9la21}.

\begin{figure} \vspace{-.43in} 
 %   \centering
  %  \includegraphics[width=0.5\linewidth]{}
\[ \hspace{-1cm} 
\left(
\begin{array}{ccccccc|ccc|ccccccc|ccc}
 \alpha  & 1 & 1 & 1 & 1 & 0 & 0 & 0 & 0 & 0 & 0 & 0 & 0 & 0 & 0 & 0 & 0 & 0 & 0 & 0 \\
 1 & \alpha  & 1 & 1 & 0 & 1 & 0 & 0 & 0 & 0 & 0 & 0 & 0 & 0 & 0 & 0 & 0 & 0 & 0 & 0 \\
 1 & 1 & \alpha  & 0 & 1 & 1 & -\frac{1}{2} & 0 & 0 & 0 & 0 & 0 & 0 & 0 & 0 & 0 & -\frac{\sqrt{3}}{2} & 0 & 0 & 0 \\
 1 & 1 & 0 & \alpha  & 1 & 1 & 0 & 0 & 0 & 0 & 0 & 0 & 0 & 0 & 0 & 0 & 0 & 0 & 0 & 0 \\
 1 & 0 & 1 & 1 & \alpha  & 1 & -\frac{1}{2} & 0 & 0 & 0 & 0 & 0 & 0 & 0 & 0 & 0 & \frac{\sqrt{3}}{2} & 0 & 0 & 0 \\
 0 & 1 & 1 & 1 & 1 & \alpha  & 1 & 0 & 0 & 0 & 0 & 0 & 0 & 0 & 0 & 0 & 0 & 0 & 0 & 0 \\
 0 & 0 & -\frac{1}{2} & 0 & -\frac{1}{2} & 1 & \alpha  & 1 & 0 & 0 & 0 & 0 & \frac{\sqrt{3}}{2} & 0 & \frac{\sqrt{3}}{2} & 0 & 0 & 0 & 0 & 0 
 \\  \hline  
 0 & 0 & 0 & 0 & 0 & 0 & 1 & \alpha  & 1 & 0 & 0 & 0 & 0 & 0 & 0 & 0 & 0 & 0 & 0 & 0 
 \\
 0 & 0 & 0 & 0 & 0 & 0 & 0 & 1 & \alpha  & 1 & 0 & 0 & 0 & 0 & 0 & 0 & 0 & 0 & 0 & 0 
 \\
 0 & 0 & 0 & 0 & 0 & 0 & 0 & 0 & 1 & \alpha  & 0 & 0 & 0 & 0 & 0 & 0 & 0 & 0 & 0 & 0 
 \\  \hline 
 0 & 0 & 0 & 0 & 0 & 0 & 0 & 0 & 0 & 0 & \alpha  & 1 & -1 & 1 & -1 & 0 & 0 & 0 & 0 & 0 \\
 0 & 0 & 0 & 0 & 0 & 0 & 0 & 0 & 0 & 0 & 1 & \alpha  & 1 & 1 & 0 & -1 & 0 & 0 & 0 & 0 \\
 0 & 0 & 0 & 0 & 0 & 0 & \frac{\sqrt{3}}{2} & 0 & 0 & 0 & -1 & 1 & \alpha  & 0 & 1 & -1 & -\frac{1}{2} & 0 & 0 & 0 \\
 0 & 0 & 0 & 0 & 0 & 0 & 0 & 0 & 0 & 0 & 1 & 1 & 0 & \alpha  & 1 & 1 & 0 & 0 & 0 & 0 \\
 0 & 0 & 0 & 0 & 0 & 0 & \frac{\sqrt{3}}{2} & 0 & 0 & 0 & -1 & 0 & 1 & 1 & \alpha  & 1 & \frac{1}{2} & 0 & 0 & 0 \\
 0 & 0 & 0 & 0 & 0 & 0 & 0 & 0 & 0 & 0 & 0 & -1 & -1 & 1 & 1 & \alpha  & 1 & 0 & 0 & 0 \\
 0 & 0 & -\frac{\sqrt{3}}{2} & 0 & \frac{\sqrt{3}}{2} & 0 & 0 & 0 & 0 & 0 & 0 & 0 & -\frac{1}{2} & 0 & \frac{1}{2} & 1 & \alpha  & 1 & 0 & 0 
 \\ \hline 
 0 & 0 & 0 & 0 & 0 & 0 & 0 & 0 & 0 & 0 & 0 & 0 & 0 & 0 & 0 & 0 & 1 & \alpha  & 1 & 0 \\
 0 & 0 & 0 & 0 & 0 & 0 & 0 & 0 & 0 & 0 & 0 & 0 & 0 & 0 & 0 & 0 & 0 & 1 & \alpha  & 1 \\
 0 & 0 & 0 & 0 & 0 & 0 & 0 & 0 & 0 & 0 & 0 & 0 & 0 & 0 & 0 & 0 & 0 & 0 & 1 & \alpha  \\
\end{array}
\right)
\]
    \caption{Gram matrix  $\Gamma(\Deltaa^{8}_{(6,(2,1))}) $ 
    %\ppm{[for what $\l$?]\bajm{[whoops]}}
    \label{fig:n9la21}}
\vspace{-.0251in} 
\end{figure}

\medskip  

\newcommand{\B}[2]{{\mathsf B}^{#1}_{#2}}

We can again obtain inductive formulae for determinants by expanding along rows corresponding to the end of a tail. We proceed as follows.

\mdef \label{par:bmn}
For $\l$ any partition, and $\bm{n}=(n_1,n_2,\dots,n_{d_{\l}})$ a $d_\l$-tuple of integers with $n_k \geq l+4$, we introduce the set 
\[
\BBB^{(\l)}_{(l,\bm{n})}:=
   \{ b_i \uu_{jk} \ | \ i= 1,\dots, d_{\l}, \ \ \uu_{jk}\in \BBB_{(l,n_i)}\}.
\]
We may regard $\BBB^{(\l)}_{(l,\bm{n})}$ as a subset of the basis $\BBB^{n}_{(n-2,\l)}$
for $\Delta^{n}_{(p,\l)}$, where $n=\max(\bm{n})$. Indeed 
\[
\BBB^{(\l)}_{(l,\bm{n}=(n,n,...,n))}   \;  = \;    \BBB^{n}_{(n-2,\l)}  
.\]
As such, we may consider the matrix of the form $\langle\_,\_\rangle$ on $ \B{(\lambda)}{(l,\bm{n})} \times \B{(\lambda)}{(l,\bm{n})}$
- denote this matrix by 
\[
\Gamma_{(l,\bm{n})}^{(\l)}   :=  
  \langle\_,\_\rangle|_{  \B{\lambda}{l,\bm{n}} \times \B{\lambda}{l,\bm{n}} }  . 
\] 
This matrix contains the head corresponding to each $b_i\in \mathcal{B}$, but in general can have tails of differing length for each such $b_i$. 

By expanding the determinant along the row at the end of any tail (assuming said tail is long enough), we can verify that the multisequence of determinants $\det(\Gamma_{l,\bm{n}}^{(\l)})$, satisfies the Chebyshev recursion \eqref{eq:Cheby0}, in each component of $\bm{n}$. More precisely, we have proven 

\begin{proposition} \label{prop:detrec} Fix $l$ and let $\lambda\vdash l+2$.
Let $\bm{n}=(n_1,\dots,n_{d_\l})$ be as before and let $e_k\in \Z^{d_{\l}}$ denote the elementary vector. Then we have
\beq
\det(\Gamma_{l,\bm{n}+2e_k}^{(\l)})=\a \det(\Gamma_{l,\bm{n}+e_k}^{(\l)})-\det(\Gamma_{l,\bm{n}}^{(\l)}).
\label{eq:detrec2}
\eq 
\qed\end{proposition}

That is to say, by (\ref{de:Chebyab}), 
\beq   \label{eq:detGamP}
\det(\Gamma_{l,\bm{n}}^{(\l)})=  \PP^{a,b}_{n_k} 
\eq 
for some $a,b$ 
which in general, may depend on the basis element $b_k$, and the $d_\l-1$ integers $\bm{n}_{\setminus k}$ where 
\beq  \label{eq:de-nminusk} 
\bm{n}_{\setminus k} := (n_1,n_2,...,n_{k-1},n_{k+1},...,n_{d_\lambda}).
\eq
Specifically 
\begin{align*}
a &= \det\left(\Gamma_{l,(n_1,...,n_{k-1},l+4,n_{k+1},...,n_{d_\lambda})}^{(\l)}\right),\\
b &= \det\left(\Gamma_{l,(n_1,...,n_{k-1},l+5,n_{k+1},...,n_{d_\lambda})}^{(\l)}\right).\end{align*}
It follows from \eqref{eq:adep} that $a$ and $b$ are monic polynomials and satisfy $\deg(b)=\deg(a)+1$. Therefore, $\PP^{a,b}$ can be written as 
\beq  \label{eq:PcP}
\PP^{a,b}
\; = \;  c \; \PP^{a',b'} 
\eq 
where $a$ and $b$ have the maximal (monic) common factor $c$, and $a'=a/c$ and $b'=b/c$. For $c$ chosen as such, the series $\PP^{a',b'}$ then has the ramping property for $n\geq l+4$ (\textit{cf.} \ref{def:chebycond}).


% commented bc it was giving an error
%\sout{
%\ppm{Observe that if $a,b$ have a} \bajm{(maximal)}\ppm{common factor $c$ then %$\PP^{a,b}$ can be written
%\beq  \label{eq:PcP}
%\PP^{a,b} %_{n_k} 
%\; = \;  c \; \PP^{a/c,b/c} %_{n_k}
%\eq 
%}
%}

\medskip \medskip 

The rest of \S\ref{ss:non1dsp} will be devoted to proving the following result, 
which shows that every one-cup determinant is determined by a Chebyshev recurrence from two in low rank (which are in turn amenable to direct computation, as we will illustrate in \S\ref{ss:examples0}):

\begin{theorem} \label{thm:main} 
%\ppm{[maybe this is a lemma and the thm is when we give it with specific polynomials? not sure yet.]}
Fix $l$. For any $\l\vdash l+2$, there exists a monic polynomial $C^{(\l)}$, and a series of polynomials $P_{n}^{(\l)}$ with the ramping property for $n\geq l+4$ (\textit{cf.} \ref{def:chebycond}), such that for $n\geq l+4$ 
%\ppm{
(and hence in the $\l$-arm of the rollet graph $\Roll_l$ \textit{cf.} \ref{pa:rollarms})    %[right?]}\bajm{[yep]}
\beq 
\Delttt{n}{n-2,\l}(\a) =   C^{(\l)}(\a) (P_{n}^{(\l)}(\a))^{d_{\l}},
\label{eq:detform}
\eq
%\ppm{[-either both sides depend on $\xx$ or neither does?]}\bajm{[both then]}
up to an overall complex scalar, 
\ppmm{where $d_\l =\dim(\SS_\l)$}.
%\ppm{[-once again here is seems that $\Delta$ means the matrix not the module. I am fine with this, but eventually we will have to be consistent. I'll replace $\Delta$ with the macro for the matrix, which is Delt and presently unpacks as $\Delt{}{}$ ?  (and, to get to the end, maybe we should create a glossary of terms for quick reference?)]}
%\bajm{[I am just using the convention we had in \eqref{eq:detnot} but happy to change]}
\end{theorem}

Firstly, we will make use of the following Lemma.

\begin{lemma}\label{lem:xi} 
Fix $l$. For each $\l\vdash l+2$, there exists a unique (up to a polynomial factor) %\ppm{[as opposed to polynomial scalar, ?]}
element $\xi_{\l} \in \Delta^{(l+4)}_{(l+2,\l)}$ with coefficients in $\C[\a]$, such that 
\beq
\uu_{ij}^* \, \xi_{\l} = 0, \text{ for } (i,j)\neq (l+3,l+4), \qquad 
    \uu_{l+3,l+4}^*\, \xi_{\l}  \ppmm{\;=\; D(\xx)} c_\l.  \label{eq:van}
\eq  
\ppmm{for some scalar $D(\xx)\in \C[\a]$.} 
Here $\uu_{ij}\in B_{(l,l+4)}$ as per \eqref{cupbasis}, $(\_)^*$ is as in \ref{de:flip2} and $c_\l \in J_{l}(l+2,l+2)$ is as in \ref{de:cl} (\textit{cf.} also \ref{ss:keybootl1}).%\bajm{[this statement is getting a bit murky]}\ppm{[-the last sentence? something to discuss!]}
\end{lemma}

\begin{proof} 
Set $\Delta=\Delta^{(l+4)}_{(l+2,\l)}$ for this proof. 
And let $\BB=\{b_1,b_2,\dots, b_{d_\l}\}$ be an orthonormal (with respect to $\langle\_,\_\rangle$) basis for $\SS_\l$, 
\ppmm{and hence for $\Deltaaa{n}{n,\l}$}, 
such that 
$b_1= \nC_\l \in \C \Sym_{l+2}$.  

For any $\xi \in \Delta$, the expression 
$\uu_{ij}^* \, \xi \in \Delta^{(l+2)}_{(l+2,\l)}$. 
The elements $\uu_{ij}^* \xi$ can thus be expressed  % written 
uniquely as a linear combination of the $b_k$. That is,  $\uu_{ij}^* \, \xi=\sum a_k b_k$ with coefficients $a_i$ determined by
\[  
a_k  =  \langle b_k, \uu_{ij}^* \, \xi \rangle  =   \langle \uu_{ij} b_k, \xi \rangle. 
\]
\ppmm{where the first equality uses the form for $\Deltaaa{n}{n,\l}$ and the second then uses the form for $\Deltaaa{n}{n-2,\l}$}.
%Thus, %writing 
Expanding $\xi_{}$ in the basis $\BBB^n_{n-2,\l}$ from \eqref{eq:basis}:
$\xi_{}=\sum_{i,j,k} [\xi_{}]_{i,j,k} \, \uu_{i,j} b_k$.  
Thus, \ppmm{comparing with (\ref{de:grammat}),}
equations \eqref{eq:van} 
%(assuming for now 
%%\ppm{[-or just put in a constant $\kappa$, say $D(\xx)$? otherwise the logic of the exposition is slightly stretched]} 
%$\uu_{l+3,l+4}^*\, \xi_{\l} = \;  \ppm{D(\a)} \; c_\l$) 
may be written in component form as
\[ 
\Gamma(\Delta) \cdot [\xi_{\l}] = \; \ppmm{D(\a)} \;  e_{l+3,l+4,1},  \tag{$\star$} \]
where $e_{l+3,l+4,1}$ is the elementary column vector 
(with entries being $0$, except in the position corresponding to the basis element $\uu_{l+3,l+4} b_1$, where it is $1$). 

%By \eqref{eq:adep}, 
%It follows
Note, for example from  \eqref{eq:adep}, 
that the matrix $
\Gamma(\Delta)$ is generically invertible (i.e. considering $\a$ as a formal parameter and working over the field of rational functions $\C(\a)$). Thus the equation $(\star)$ has a unique solution 
for any factor \ppmm{$D(\a)$}, 
with coefficients $[\xi_{\l}]_{i,j,k}\in \C(\a)$. 
Multiplying $\xi_\l$ by a common denominator  %, say $D(\a)$, 
of the coefficients, if necessary, we 
simply vary the factor. Hence we 
obtain the desired $\xi_\l$, obeying $\uu_{l+3,l+4}^* \, \xi_l=D(\a) c_\l$. 
\end{proof}

%\ppm{[is the rest actually a post-proof remark?]}\bajm{[yes]}

\mdef 
\ppmm{Passing to the ground field $\C$ by specialising $\xx$ to}
%\sout{Regarding $\a$ as} 
a complex number then, we note that it may be the case that $\xi_\l$ also satisfies $\uu_{l+3,l+4}^* \, \xi_l=0$ (\textit{i.e.} when $\a$ was a root of the denominator $D(\a)$).
%%\ppm{[-surely then it -is-, not maybe? nope. misreading.]}
%\end{proof}

\mdef\label{niceelt} 
Let $\xi_{\l} \in \Delta^{(l+4)}_{(l+2,\l)}$ be as per Lemma~\ref{lem:xi} 
\ppmm{ - i.e. in the first one-cup standard module in the $\l$-arm.} %\bajm{[$\l$ arm is undefined... can add this]}.
And let
$\BB=\{b_1,b_2,\dots, b_{d_{\l}}\}$ be an orthonormal basis for $\SS_\l$, with $b_k=x_k \nC_\l$ and $x_k\in \C S_{l+2}$.
For $k =1,2,...,d_\l$ 
define 
\beq 
\xi_{\l}^{(b_k)}:= (x_k \nC_\l \otimes \text{id}_2) \, \xi_\l \in \Delta^{(l+4)}_{(l+2,\l)}, 
\eq
regarding $\nC_\l$ as an element of $\C \Sym_{l+2}\subset J_{l,l+2}$. We now claim the $\xi_{\l}^{(b_k)}$ obey the following: 

\begin{lemma}\label{lem:niceelt} Fix $l$, let $\l \vdash l+2$. Let $\BB$ and $\xi_{\l}^{(b_k)} \in \Delta^{(l+4)}_{(l+2,\l)} $ be as per \ref{niceelt}. Then $\xi_{\l}^{(b_k)}$ has coefficients in $\C[\a]$ and satisfies:
\begin{enumerate}
    \item The coefficient of the basis vector $\uu_{l+3,l+4} \, b_{k'}$ in $\xi_{\l}^{(b_k)}$ is zero unless $k'=k$. Furthermore, the coefficient of $\uu_{l+3,l+4} \, b_{k}$ is independent of $b_k$.
    \item  The form $\langle \xi_{\l}^{(b_k)}, \uu_{i,j} b_{k'}  \rangle $ vanishes unless $(i,j)=(l+3,l+4)$ and $k'=k$. Furthermore, $\langle \xi_{\l}^{(b_k)}, \uu_{l+3,l+4} b_{k}  \rangle $ is independent of $k$.
\end{enumerate}
\end{lemma}
\begin{proof} 
Firstly, since $\xi_\l$ had coefficients in $\C[\a]$ and since $x_i, \nC_\l \in \C \Sym_{l+2}$, it follows that $\xi_\l^{(b_k)}$ has coefficients in $\C[\a]$. 

For any cup $\uu_{i,j}$ with $j<l+4$, the composite $(x_i \nC_\l  \otimes \text{id}_2)\uu_{i,j}$ can be written as a combination of the form $\sum_{i'<j'<l+4} \uu_{i',j'} x_{i',j'}$ with $x_{i',j'}\in \C\Sym_{l+2}$, by ``pulling" the cup up through each permutation. Therefore, the only contribution to basis vectors of the form  $\uu_{l+3,l+4} b_{k''}$ in $\xi_\l^{(b_k)}$ comes from basis vectors of the form $\uu_{l+3,l+4} b_{k'}$ in $\xi_\l$. Computing
\begin{align*}(x_k \nC_\l  \otimes \text{id}_2) \uu_{l+3,l+4} b_{k'}&=\uu_{l+3,l+4}(x_k \nC_\l x_{k'} c_{\l})\\
&=\langle e c_\l, b_{k'}\rangle \  \uu_{l+3,l+4}x_k \nC_{\l} c_{\l}\\
&=\langle e c_{\l}, b_{k'}\rangle \ \uu_{l+3,l+4}b_k,\end{align*}
(using the fact that $\nC_{\l} c_{\l}=c_{\l}$) it follows that the coefficient of $\uu_{l+3,l+4} \, b_{k''}$ (in $\xi_\l^{(b_k)}$) is 0 unless $k''=k$, in which case we have
\[ \left[\xi_\l^{(b_k)}\right]_{l+3,l+4,k}= \sum_{k'} \left[\xi_\l \right]_{l+3,l+4,k'} \langle e c_\l ,b_{k'}\rangle,\]
which is manifestly independent of $k$, proving part 1 of the Lemma.

To prove part 2, we first note that for any $i<j<l+4$, $\uu_{i,j}^* (x_k \nC_\l\otimes \text{id}_2)$ can be written as a combination $\sum_{i'<j'<l+4} x_{i',j'}\uu^*_{i',j'}$ with  $x_{i',j'}\in \C \Sym_{l+2}$ (by ``pulling" a cap down through each permutation). Therefore, by \eqref{eq:van} it follows that for any such $i<j$ we have $\uu_{i,j}^*\xi_{\l}^{(b_k)}=0$ and in particular $\langle \uu_{i,j} b_{k'}, \xi_{\l}^{(b_k)}   \rangle =0$ for any $k'$. When $(i,j)=(l+3,l+4)$, we have
\begin{align*}
    c_{\l} x_{k'}^* \uu_{l+3,l+4}^* \xi_{\l}^{(b_k)}&= 
    c_{\l} x_{k'}^* \uu_{l+3,l+4}^* (x_k \nC_\l \otimes \text{id}_2) \, \xi_\l,\\
    &=(c_\l \,  x_{k'}^* \, x_k \, \nC_\l) \uu_{l+3,l+4}^* \, \xi_\l\\
    &=\langle b_{k'} , b_k\rangle \nC_\l \uu_{l+3,l+4}^* \, \xi_\l=\delta_{k,k'}D(\a) c_\l, 
\end{align*}
where we have used orthogonality of the $b_k$, $\nC_\l c_\l=c_\l$, and $\uu_{l+3,l+4}^* \, \xi_\l= D(\a) c_\l$ (as per \ref{lem:xi}) in the last line. We can thus read off that 
\[
\langle \uu_{l+3,l+4} b_{k'}, \xi_{\l}^{(b_k)}  \rangle =\delta_{k,k'} D(\a),
\]
which completes the proof since $D(\a)\in \C[\a]$ is manifestly independent of $k$. 
\end{proof}

We now apply this Lemma to prove the following:

\begin{lemma}\label{lem:pl} Fix $l$, let $\l \vdash l+2$. Let $\BB$ be an orthonormal basis for $\SS_\l$, let $\bm{n}=(n_1,\dots, n_{d_\l})$ be a $d_\l$-tuple of integers with $n_k\geq l+4$ for all $k$, and let $\Gamma_{(l,\bm{n})}^{(\l)}$ be as per \ref{par:bmn}. Then for each $k=1,\dots, d_\l$, there exists a monic polynomial $C^{(\l,\bm{n}_{\setminus k})}$ and a series of polynomials $P^{(\l)}_{n}$ with the ramping property for $n\geq l+4$ (\textit{cf.} \ref{def:chebycond}), such that 
\[ 
\text{det}(\Gamma_{(l,\bm{n})}^{(\l)}) \;=\; C^{(\l,\bm{n}_{\setminus k})}(\a) \; P^{(\l)}_{n_k}(\a),
\]
where $\bm{n}_{\setminus k}$ is as per \eqref{eq:de-nminusk}. In particular, $P^{(\l)}_{n_k}$ (and hence $C^{(\l,\bm{n}_{\setminus k})}$) is independent of the basis vector $b_k\in \BB$. Furthermore, $P^{(\l)}_{n_k}$ does not depend on the integers $\bm{n}_{\setminus k}$.
\end{lemma}

\begin{proof} 
Monicity of the determinant $\text{det}(\Gamma_{(l,\bm{n})}^{(\l)})$, its degree, and the recursion established in Prop.\ref{prop:detrec}, imply that for each such $k$, we may write   
    \[
    \text{det}(\Gamma_{(l,\bm{n})}^{(\l)})=C^{(b_k,\l,\bm{n}_{\setminus k})}(\a) P^{(b_k,\l, \bm{n}_{\setminus k})}_{n_k}(\a),
    \]
where $C^{(b_k,\l,\bm{n}_{\setminus k})}(\a)$ is a monic polynomial and $P^{(b_k,\l, \bm{n}_{\setminus k})}_{n}(\a)$ is a series of polynomials with the ramping property for $n\geq l+4$. In general, $P^{(b_k, \l,\bm{n}_{\setminus k})}_{n_k}(\a)$ may depend on $b_k$ and on the integers $\bm{n}_{\setminus k}$; the crux of this proof is thus to show otherwise. To do this, it is sufficient to verify that for any $k=1,\dots, d_\l$ and choice of integers $\bm{n}_{\setminus k}$, the sequence has the same ``initial conditions", that is,
\[
P^{(b_k,\bm{n}_{\setminus k})}_{l+4}(\a)=a(\a), \qquad P^{(b_k,\bm{n}_{\setminus k})}_{l+5}(\a)=b(\a),  
\]
for polynomials $a,b$ independent of $b_k$ and $\bm{n}_{\setminus k}$. Furthermore, by \ref{def:chebycond} it suffices to demonstrate the ratio 
\beq
\frac{P^{(b_k,\bm{n}_{\setminus k})}_{l+4}(\a)}{P^{(b_k,\bm{n}_{\setminus k})}_{l+5}(\a)}=\frac{a(\a)}{b(\a)}=\frac{\det(\Gamma^{(\l)}_{(l,(n_1,\dots, n_{k-1},l+4,n_{k+1},\dots, n_{d_\l}))})}{\det(\Gamma^{(\l)}_{(l,(n_1,\dots, n_{k-1},l+5,n_{k+1},\dots, n_{d_\l}))})},\label{eq:ratio}
\eq
is independent of $b_k$ and $\bm{n}_{\setminus k}$.

Suppose, by reordering $\BB$ if necessary, that $k=d_\l$. Now for a $(d_\l-1)$-tuple $\bm{n'}=(n'_1,\dots,n'_{d_\l-1})$, with $n'_{k'}\geq l+4$ for all $k'$, consider the matrix $\Gamma:=\Gamma^{(\l)}_{(l,(\bm{n'},l+5))}$. By the computation \eqref{cheby}, it has the form 
\beq
\Gamma=\Gamma^{(\l)}_{(l,(\bm{n'},l+5))}=
\begin{pNiceArray}{cccc|c}
\Block{4-4}<\Large>{\Gamma^{(\l)}_{(l,(\bm{n'},l+4))}} &&&& 0\\
&&&& \vdots\\
&&&& 0 \\
&&&& 1 \\ \hline
0 & \dots &0 &1 & \a \\
\end{pNiceArray}.\label{eq:matlastrow}
\eq

By Lemma \ref{lem:niceelt} part 2, it follows that the following linear combination of the rows of the Gram-matrix $\Gamma(\Delta^{(l+4)}_{(l+2,\l)} )=\Gamma^{(\l)}_{(l,\bm{l'})}$ (where $\bm{l'}$ is the $d_\l$-tuple with all entries being $l+4$) yields the row vector
\begin{align*}
    \sum_{i,j,k} \left[\xi^{(b_{d_\l})}_{\l}\right]_{i,j,k} R_{\uu_{i,j} b_k}&=  \langle \xi^{(b_{d_\l})}_{\l}, \uu_{l+3,l+4} \, b_{d_{\l}}\rangle \, e_{l+3,l+4, d_{\l}}\\
    &=\begin{pmatrix}
    0 &\dots & 0 & \langle \xi^{(b_{d_\l})}_{\l}, \uu_{l+3,l+4} \, b_{d_{\l}}\rangle
\end{pmatrix}
\end{align*}
where $e_{i,j,k}$ denotes the elementary vector. Furthermore, by \ref{lem:niceelt} part 1 this combination does not include the rows $R_{\uu_{l+3,l+4} b_{k}}$ for $k<d_{\l}$. Since the row $R'_{\uu_{i,j} b_k}$ with $j<l+4$ in the matrix $\Gamma$ can be obtained from the corresponding row $R_{\uu_{i,j} b_k}$ in the Gram-matrix $\Gamma(\Delta^{(l+4)}_{(l+2,\l)})$, by putting a zero in any column labelled by an element of any tail (that is, an element of the form $\uu_{i,j} b_k$ with $j>l+4$), it follows that the same combination of the corresponding rows of the matrix $\Gamma$ yields the row vector
\[ \begin{pmatrix}
    0 &\dots & 0 & \langle \xi^{(b_{d_\l})}_{\l}, \uu_{l+3,l+4} b_{d_{\l}}\rangle & \left[\xi^{(b_{d_\l})}_{\l}\right]_{l+3,l+4,d_{\l}}
\end{pmatrix}.\]
Therefore, after multipling the last row of $\Gamma$ by $\langle \xi^{(b_{d_\l})}_{\l}, \uu_{l+3,l+4} b_{d_{\l}}\rangle$, we can apply a sequence  of row reductions to exhibit the equality
\begin{align*}
    \langle \xi^{(b_{d_\l})}_{\l}, \uu_{l+3,l+4} b_{d_{\l}}\rangle \det(\Gamma) &=
\begin{vNiceArray}{cccc|c}
\Block{4-4}<\Large>{\Gamma^{(\l)}_{(l,(\bm{n'},l+4))}} &&&& 0\\
&&&& \vdots\\
&&&& 0 \\
&&&& 1 \\ \hline
0 & \dots &0 & \langle \xi^{(b_{d_\l})}_{\l}, \uu_{l+3,l+4} b_{d_{\l}}\rangle & \a \, \langle \xi^{(b_{d_\l})}_{\l}, \uu_{l+3,l+4} b_{d_{\l}}\rangle\\
\end{vNiceArray}\\
&=\begin{vNiceArray}{cccc|c}
\Block{4-4}<\Large>{\Gamma^{(\l)}_{(l,(\bm{n'},l+4))}} &&&& 0\\
&&&& \vdots\\
&&&& 0 \\
&&&& 1 \\ \hline
0 & \dots &0 & 0 & \a \, \langle \xi^{(b_{d_\l})}_{\l}, \uu_{l+3,l+4} b_{d_{\l}}\rangle-\left[\xi^{(b_{d_\l})}_{\l}\right]_{l+3,l+4,d_{\l}}\\
\end{vNiceArray}\\
&=\left( \a \, \langle \xi^{(b_{d_\l})}_{\l}, \uu_{l+3,l+4} b_{d_{\l}}\rangle-\left[\xi^{(b_{d_\l})}_{\l}\right]_{l+3,l+4,d_{\l}} \right)\det(\Gamma^{(\l)}_{(l,(\bm{n'},l+4))}).
\end{align*}
Comparing with \eqref{eq:ratio}, we now obtain
\[\frac{P^{(b_{d_{\l}},\bm{n'})}_{l+4}(\a)}{P^{(b_{d_\l},\bm{n'})}_{l+5}(\a)}=\frac{\langle \xi^{(b_{d_\l})}_{\l}, \uu_{l+3,l+4} b_{d_{\l}}\rangle}{\left( \a \, \langle \xi^{(b_{d_\l})}_{\l}, \uu_{l+3,l+4} b_{d_{\l}}\rangle-\left[\xi^{(b_{d_\l})}_{\l}\right]_{l+3,l+4,d_{\l}} \right)},\]
which is manifestly independent of $\bm{n'}$, and independent of $b_{d_{\l}}$ by Lemma \ref{lem:niceelt}. This completes the proof. \end{proof}

We now apply this Lemma to prove:

\begin{proposition}\label{prop:main}
    Fix $l$, let $\l \vdash l+2$. Let $\BB$ be an orthonormal basis for $\SS_\l$, let $\bm{n}=(n_1,\dots, n_{d_\l})$ be a $d_\l$-tuple of integers $n_k\geq l+4$ and let $\Gamma_{(l,\bm{n})}^{(\l)}$ be as per \ref{par:bmn}. There exists a monic polynomial $C^{(\l)}$ and a series of polynomials $P^{(\l)}_{n}$ with the ramping property for $n\geq l+4$ (\textit{c.f} \ref{def:chebycond}), such that 
\[ \text{det}(\Gamma_{(l,\bm{n})}^{(\l)})=C^{(\l)}(\a) \prod_{k=1}^{d_{\l}} P^{(\l)}_{n_k}(\a).\]\end{proposition}
\begin{proof} By Lemma \ref{lem:pl}, it follows that each $P^{(\l)}_{n_k}(\a)$ divides $\text{det}(\Gamma_{(l,\bm{n})}^{(\l)})$. It suffices to prove that $\prod_{k=1}^{d_{\l}} P^{(\l)}_{n_k}(\a)$ divides $\text{det}(\Gamma_{(l,\bm{n})}^{(\l)})$ since then, comparing with the factorisation from Lemma \ref{lem:pl} yields
\[ C^{(\l)}(\a) \prod_{k\neq k' } P^{(\l)}_{n_k}(\a)= C^{(\l,\bm{n}_{\setminus k'})}(\a), \]
for each $k'=1,\dots, d_{\l}$ which shows $C^{(\l)}(\a)$ is independent of all the integers $n_{k'}$. 

Suppose for a contradiction that for some $\bm{n}$, $\prod_{k=1}^{d_{\l}} P^{(\l)}_{n_k}(\a)$ does not divide $\text{det}(\Gamma_{(l,\bm{n})}^{(\l)})$. By reordering the basis $\BB$ is necessary, it follows there is some integer $2\leq d \leq d_{\l}$ such that $\prod_{k=1}^{d-1} P^{(\l)}_{n_k}(\a)$ divides $\text{det}(\Gamma_{(l,\bm{n'})}^{(\l)})$ for any $d_\l$-tuple $\bm{n'}$ which agrees with $\bm{n}$ in the first $(d-1)$-entries, yet $\prod_{k=1}^{d} P^{(\l)}_{n_k}(\a)$ does not divide $\text{det}(\Gamma_{(l,\bm{n})}^{(\l)})$. Since $\prod_{k=1}^{d-1} P^{(\l)}_{n_k}(\a)$ divides $\text{det}(\Gamma_{(l,\bm{n})}^{(\l)})$ by assumption and since $P^{(\l)}_{n_{d}}(\a)$ must divide it, it follows these two factors must ``overlap" non-trivially. That is, we must be able to write
\[ \text{det}(\Gamma_{(l,\bm{n})}^{(\l)})= A(\a) B(\a) C(\a) D(\a),\]
where $D(\a)$ has positive degree and does not divide $A(\a)$, such that $B(\a) D(\a)= P^{(\l)}_{n_{d}}(\a)$ and $C(\a) D(\a)= \prod_{k=1}^{d-1} P^{(\l)}_{n_k}(\a)$. But then lemma \ref{lem:pl} implies that   
\[\text{det}(\Gamma_{(l,\bm{n}+e_{d})}^{(\l)})= A(\a) C(\a)  P^{(\l)}_{n_{d}+1}(\a). \]
Since $C(\a)D(\a)$ must divide $\text{det}(\Gamma_{(l,\bm{n}+e_{d})}^{(\l)})$, yet $D$ does not divide $A$, it follows that $D$ must divide $P^{(\l)}_{n_{d}+1}$. The Chebyshev recursion then implies that $D$ is a non-trivial common factor of the sequence $P^{(\l)}_{n}$ which is a contradiction.\end{proof}

We are now able to prove Theorem \ref{thm:main} as a specialisation of Prop. \ref{prop:main}:

\begin{proof}[Proof of Theorem \ref{thm:main}] For any $n\geq l+4$, taking $\bm{n}$ to be the $d_{\l}$-tuple, $\bm{n}=(n,n,\dots, n)$, yields
\[\Delttt{(n)}{n-2,\l}=\text{det}(\Gamma_{(l,\bm{n})}^{(\l)})=C^{(\l)}(\a) (P^{(\l)}_{n}(\a))^{d_{\l}},\]
as desired. \end{proof}


\subsection{Applications of Theorem~\ref{thm:main}: \ppmm{polynomials,} and examples}   \label{ss:examples0}

In this subsection we apply Theorem \ref{thm:main} to obtain closed forms for gram determinants of the $J_{l,n}$-modules, $\Delt{n}{n-2,\l}$ with $n\geq l+4$, for all cases when $l\leq 2$, by direct computation of low rank (namely $n=l+4,l+5$) cases. 
%\ppm{[-explain where this restriction comes from etc.]}\bajm{[it comes from us??]}\ppm{[yes. somewhere above!]}
We also study the cases $\l=(l+2), (1^{l+2}), (l+1,1)\vdash l+2$ and $\l=(l+1,1)^T\vdash l+2$ for $l>0$. We begin by giving examples for $\l$ such that $\SS_\l$ is 1-dimensional.

\mdef \label{pa:1dgdets}
For $l=0$  %, we consider the row and column partitions $\l$ here.  In these cases, 
direct computation yields
\begin{align*}
\Delttt{(4)}{2,(2)}&=(\alpha -1) \alpha  \left(\alpha ^2+\alpha -4\right),& \Delttt{(5)}{3,(2)}&=(\alpha -1) \left(\alpha ^4+\alpha ^3-5 \alpha ^2-\alpha +2\right),\\
    \Delttt{(4)}{2,(1^2)}&=(\alpha -2) (\alpha -1) (\alpha +1) (\alpha +2),&
   \Delttt{(5)}{3,(1^2)}&=(\alpha -1) (\alpha +2) \left(\alpha ^3-\alpha ^2-3 \alpha +1\right),
\end{align*}  
For $l=1$, we consider the row and column partitions $\l$ here. 
In these cases, direct computation yields
\begin{align*}
%   \\
    \Delttt{(5)}{3,(3)}&=(\alpha -2)^2 \alpha ^2 (\alpha +1) \left(\alpha ^2+3 \alpha -6\right),& \Delttt{(6)}{4,(3)}&=(\alpha -2)^2 \alpha ^3 \left(\alpha ^3+4 \alpha ^2-4 \alpha -10\right),\\
    \Delttt{(5)}{3,(1^3)}&=(\alpha -3) (\alpha -2)^2 (\alpha +1) (\alpha +2)^3,&
    \Delttt{(6)}{4,(1^3)}&=(\alpha -2)^2 (\alpha +2)^3 \left(\alpha ^3-2 \alpha ^2-4 \alpha +2\right).
\end{align*}
Thus, using the notation of Theorem \ref{thm:main}, 
%\ppm{[-at the mo this is a forward ref!]}\bajm{[no more]}
we have here $\Delttt{(n)}{n-2,\l}=C^{(\l)}(\a) P^{(\l)}_{n}(\a)$ where
\begin{align*}
    C^{(2)}(\a)&=(\a-1), &    C^{(1^2)}(\a)&=(\a-1)(\a+2),\\
    C^{(3)}(\a)&=(\a-2)^2\a^2, &    C^{(1^3)}(\a)&=(\a-2)^2(\a+2)^3,
\end{align*}
and the sequences $P_n^{(\l)}(\a)$ are completely determined by the Chebyshev recursion, and the ``initial conditions":
\begin{align*}
    P_4^{(2)}(\a)&=\a(\a^2+\a-4), &        P_5^{(2)}(\a)&=\a^4+\a^3-5\a^2-\a+2,\\
    P_4^{(1^2)}(\a)&=(\a-2)(\a+1), &        P_5^{(1^2)}(\a)&=\alpha ^3-\alpha ^2-3 \alpha +1,\\
    P_5^{(3)}(\a)&=(\alpha +1) \left(\alpha ^2+3 \alpha -6\right), &        P_6^{(3)}(\a)&=\alpha  \left(\alpha ^3+4 \alpha ^2-4 \alpha -10\right),\\
    P_5^{(1^3)}(\a)&=(\alpha +1) \left(\a-3\right), &        P_6^{(1^3)}(\a)&=  \left(\alpha ^3-2 \alpha ^2-4 \alpha +2\right).
\end{align*}
By making use of \ref{pa:chebslv} we may write each polynomial $P_n^{(\l)}$ explicitly. We address this in this next paragraph, where we treat more cases with a general formula.


\mdef   \label{pa:able} \bajmm{By direct computation of Gram matrices $\Delttt{(n)}{n-2,\l}$ for $\l=(l+2)$ and $(1^{l+2})$, when $-1 \leq l \leq 12$ and $n=l+4,l+5$, we have 
(using the notation of Theorem \ref{thm:main})  }
%We do this in more generality; we are able to compute 
%\ppm{[-let's reword to show the proof.]}\bajm{[what this means?]}\ppm{[ah I see, more generality does not mean generality. ok I'll just think of a wording a bit different from ``are able to''.]}
%\ppm{[-forward ref.]}
%for $0\leq l\leq 12$, that
\begin{align}
    C^{(l'+2)}(\a)&=(\a+(l'-1))^{l+1}(\a-2)^{\frac{1}{2}l'(l'+3)},\label{eq:crow}\\
    C^{(1^{l+2})}(\a)&=(\a-(l+1))^{l+1}(\a+2)^{\frac{1}{2}(l+2)(l+1)},\label{eq:ccol}
\end{align}
and
\begin{align}P^{(l'+2)}_{l'+4+k}&=P_{k+4}^U(\a)+(3 l'+1) P_{k+3}^U(\a)+
\left(2 l'^2-l'-2\right)P_{k+2}^U(\a)\nn\\
&\quad+(l'+1)((2l'-1)P^U_{k+1}(\a)-P^U_{k}(\a)),\label{eq:prow}\\
P^{(1^{l+2})}_{l+4+k}&=P_{k+3}^U(\a)-(l+1)(P_{k+2}^U(\a)+P_{k+1}^U(\a)).\label{eq:pcol}\end{align}
These formulas are verified for $-1\leq l\leq 12$ and $0\leq l'\leq 12$.

\begin{conjecture} Formulae \eqref{eq:crow},\eqref{eq:ccol}, \eqref{eq:prow} and \eqref{eq:pcol} hold for all $l\geq -1$ and $l'\geq 0$. 
\end{conjecture}

\mdef By direct computation, using formulae \eqref{eq:prow} and \eqref{eq:pcol} (including $l,l'>12$) we have:
\begin{align}
    P^{(l'+2)}_{l'+4}&=(\alpha +l') \left(\alpha ^2+(2l'+1)\alpha -2(l'+1)\right), \label{eq:row0}\\
    P^{(1^{l+2})}_{l+4}&=(\alpha +1) (\alpha- (l+2) ). \label{eq:col0}
\end{align}


\mdef \label{pa:21dets}
 As we have seen in \S\ref{ss:sdmods}, for $l=1$ we have gram determinants 
\begin{align*}
    \Delttt{(5)}{3,(2,1) }&= (\alpha -2)^3 \alpha  (\alpha +2) (\alpha +4)(\alpha ^4-7 \alpha ^2+3)^2,\\
    \Delttt{(6)}{4,(2,1) }&= (\alpha -2)^3 \alpha  (\alpha +2) (\alpha +4)(\alpha -1)^2 \alpha^2(\alpha +1)^2 \left(\alpha ^2-7\right)^2.
\end{align*} 
Comparing with \eqref{eq:detform} from the Theorem, \bajmm{for $n\geq 5$,} 
we have
\[
\Delttt{(n)}{n-2,(2,1) }= C^{(2,1)}(\a) (P^{(2,1)}_n(\a))^2, \quad 
%      {{C}}^{(2,1)}(\a) {{= (\alpha -2)^3 \alpha  (\alpha +2) (\alpha +4),}}
\]
%\ppm{[-seems key to say some words here about initial conditions, $n\geq l+4$ and so on, no? in the Thm $n$ is the overall rank, not the distance down the arm, as it were.]}
Since $P^\l$ is reduced we deduce
\[
C^{(2,1)}(\a)= (\alpha -2)^3 \alpha  (\alpha +2) (\alpha +4),  
\]
as the largest common factor of the two determinants. Hence we deduce
%\ppm{and hence} %where 
$P^{(2,1)}_n=\PP^{a,b}_n$ with initial conditions (at $n=5=l+4$ and $n=6=l+5$)
%\ppm{[-presumably at $n=5,6$? so maybe need a quick comment here.]}
\[
a=\alpha ^4-7 \alpha ^2+3, \qquad b=(\alpha -1) \alpha  (\alpha +1) \left(\alpha ^2-7\right).
\]
Using \ref{pa:chebslv} we can determine
\begin{align*}
    P^{(2,1)}_n&=P^U_{n}(\a)-4P^U_{n-2}(\a)-4P^U_{n-4}(\a)-2 P^U_{n-6}(\a).
\end{align*}

\mdef 
Similarly, for $l=2$, direct computation \ppmm{of the low rank gram determinants and applications of the Theorem \ref{thm:main}} yields
\begin{align*}
    C^{(3,1)}(\a)&=(\alpha -3)^5 (\alpha -2)^2 (\alpha -1)^3 (\alpha +1)^6 (\alpha +4) (\alpha +6),\\
    C^{(2,2)}(\a)&=(\alpha -3)^3 (\alpha -1)^6 (\alpha +1)^3 (\alpha +4)^4,\\
    C^{(2,1,1)}(\a)&=(\alpha -3)^6 (\alpha -1)^3 (\alpha +1)^5 (\alpha +2)^3 (\alpha +4)^4,\end{align*}
and
\begin{align*}
    P^{(3,1)}_{n}(\a)&=P^U_{n}(\a)+2 P^U_{n-1}(\a)-9 P^U_{n-2}(\a)\\
    &\quad-4 P^U_{n-5}(\a)- P^U_{n-6}(\a)-2 P^U_{n-7}(\a)-3 P^U_{n-8},\\
    P^{(2,2)}_{n}(\a)&=P^U_{n-2}(\a)-2 P^U_{n-3}(\a)-5 P^U_{n-4}(\a)+6 P^U_{n-5}(\a)-3 P^U_{n-6}(\a), \\
    P^{(2,1,1)}_{n}(\a)&=P^U_{n-1}(\a)-2 P^U_{n-2}(\a)-6 P^U_{n-3}(\a)\\
    &\quad +6 P^U_{n-4}(\a) -2 P^U_{n-5}(\a)+2 P^U_{n-6}(\a)-3 P^U_{n-7}(\a).
\end{align*}

\mdef For $2<l\leq 7$, direct computation of the Gram determinants $\Delttt{(n)}{n-2,\l}$ with $\l=(l+1,1), (l+1,1)^T$ and $n=l+4,l=5$, and application of Theorem \ref{thm:main} yields
%\ppm{[-we can try to express this more progressively.]}\bajm{[I don't know what this means]}\ppm{[same as in \ref{pa:able}.]}
\begin{align}
    C^{(4,1)}(\a)&=(\alpha -3)^{16} (\alpha -2)^5 \alpha ^{11} (\alpha +2)^{10} (\alpha +6) (\alpha +8)\nn\\
    C^{(l'+1,1)}(\a)&=(\alpha -3)^{(l'-1)(l+1)(l+3)/3} \, (\alpha -2)^{(l'-1)(l+2)/2} \, \alpha ^{l'('l+3)/2}(\alpha +(l'-3))^{l'(l'+1)/2} \nn\\
    &\quad\times (\alpha +(l'-1))^{(l'+1)(l'+2)/2} (\alpha +2l') (\alpha +2(l'+1)),\label{eq:chook}
\end{align}
where the last formula is valid for $l'\geq 4$, and
\begin{align}
        C^{(l+1,1)^T}(\a)&=(\a+4)^{l(l+1)(l+2)/6}(\a+2)^{l(l+1)/2}(\a+1)^{(l-1)(l+1)(l+3)/3}\nn\\
    &\quad\times(\alpha -(l-1))^{l(l+1)/2} (\alpha -(l+1))^{(l+1)(l+2)/2},\label{eq:chookt}
\end{align}
which holds for $l\geq 1$. The corresponding Chebyshev series are given by
\begin{align}
    P_{l''+4+k}^{(l''+1,1)}&=P^U_{k+6}(\a)+2 (2 l''-3) P^U_{k+5}(\a)+(5 l''^2-21 l''+13)P^U_{k+4}(\a)\nn\\
    &\quad +(l''-2) \left[\left(2 l''^2-17 l''+7\right)P^U_{k+3}(\a)- \left(6 l''^2-19 l''+5\right) P^U_{k+2}(\a)\right]\nn\\
    &\quad +4 \left(l''^3-6 l''^2+8 l''-1\right)P^U_{k+1}(\a) -(l''-3) \left(2 l''^2-4 l''-1\right)P^U_{k}(\a)\nn\\
    &\quad-(3 l''^2-3 l''-4)P^U_{k-1}(\a)-(l''+1)P^U_{k-2}(\a),\label{eq:phook}\end{align}
which holds for $l''\geq2$, and 
\begin{align}
    P_{l+4+k}^{(l+1,1)^T}&=P^U_{k+5}(\a)-2 (l-1) P^U_{k+4}(\a)+l(l-5)P^U_{k+3}(\a)+3l(l-1)P^U_{k+2}(\a)\nn\\
    &\quad 2(l^2-2l-1) P^U_{k+1}(\a)+l(l-1)P^U_{k}(\a)-(l+1)P^U_{k-1}(\a),\label{eq:phookt}
\end{align}
which holds for $l\geq 1$.

 \begin{conjecture}Formulae \eqref{eq:chook},\eqref{eq:chookt},\eqref{eq:phook}, and \eqref{eq:phookt}, hold for all $l\geq 1$, $l'\geq 4$ and $l''\geq 2$.     
 \end{conjecture}
\medskip 

Let us note, therefore, that in all cases directly computed, Theorem \ref{thm:main} can be strengthened as follows.  %\ppm{[-again we can maybe express more indicatively if poss.]}

\begin{proposition}\label{prop:ints} For $\l \vdash l+2$, such that either, $l\leq 2$; $\l=(l+2), (1^{l+2})$ and $-1\leq l\leq 12$; $\l=(l+1,1), (l+1,1)^T$ and $1\leq l \leq 7$, the polynomials $C^{(\l)}$ and $P_n^{(\l)}$ as per Theorem \ref{thm:main} have integer coefficients. Furthermore, each such $C^{(\l)}$ has only integer roots. \end{proposition}

%\mdef 
%\ppm{[first state the prop that is the conjecture where known true?]}


\begin{conjecture} Proposition \ref{prop:ints} holds for all $l\geq -1$ and $\l\vdash l+2$.
\end{conjecture}

\begin{remark}\label{rem:polys} 
For a given $l$ and $\lambda\vdash l+2$, we note that whilst the polynomials $P_n^{(\l)}$ only appear as a factors in a Gram-determinant for $n\geq l+4$, the same polynomial
series 
with $n<l+4$ are uniquely determined by the Chebyshev recursion. For the cases we have computed, we give $P_{l+3}^{(\l)}$ and $P_{l+2}^{(\l)}$ below:
\begin{align*}
    P_{5}^{(2,2)}(\a)&=(\a-3)(\a+1),  &P_4^{(2,2)}(\a)&=4 (\alpha -2), \\
    P^{(l+2)^T}_{l+3}(\a)&=\alpha -(l+1), &  P^{(l+2)^T}_{l+2}(\a)&=l+2,\\
    P^{(l'+2)}_{l'+3}(\a)&=(\alpha +l'-1) (\alpha +2 (l'+1) ), &  P^{(l'+2)}_{l'+2}(\a)&=(l'+2) (\alpha +2 l'),\\
    P^{(l''+1,1)^T}_{l''+3}(\a)&=(\alpha +2) (\alpha -(l''-1) ) (\alpha -(l''+1) ), &  P^{(l''+1,1)^T}_{l''+2}(\a)&=(l''+2)(\alpha +1) (\alpha -l'' ),\end{align*}
    where these formulae are valid for heights $l\geq -1$, $l'\geq 0$, and $l''\geq 1$. Furthermore, when $l\geq 2$ we have
\begin{align*}
    P^{(l+1,1)}_{l+3}(\a)&=(\alpha -2) (\alpha +l-3) (\alpha +l-1) (\alpha +2 l), \\
     P^{(l+1,1)}_{l+2}(\a)&=(l+2) (\alpha +l-2) \left(\alpha ^2+(2l-3) \alpha -2 l\right).
\end{align*}
 Remarkably then, we observe that for all cases computed thus far, $P_{l+3}^{(\l)}(\a)$ is monic, has only integer roots, and (up to a scalar factor) divides the King polynomial (see \cite[\S 5.1]{DeVisscherMartin2016}).  Furthermore, in all such cases $P_{l+3}^{(\l)}(\a)$ has an overall integer factor of $l+2$, and comparing with \eqref{eq:row0} and \eqref{eq:col0} we observe that, for $l\geq 1$
\begin{align*}
    P^{(l+1,1)^T}_{l+2}(\a)&=(l+2) P^{(l)^T}_{l+2}(\a), & P^{(l+1,1)}_{l+2}&=(l+2)  P^{(l)}_{l+2}(\a). 
\end{align*}
We also note the equalities $P^{(2,2)}_{5}(\a)=P^{(1^3)}_{5}(\a)$ and $P^{(2,2)}_{4}(\a)=4 P^{(1^3)}_{4}(\a)$.
\end{remark}



%\subsub
\section{Roots of the polynomials $P^{(\l)}_n$}\label{s:roots}
In this section, we make some observations about the distribution of the roots of the polynomials $P^{(\l)}_n$ in the cases considered. We begin by noting that all roots of these polynomials must be real; since the form $\langle \_,\_ \rangle$ on $\Delta^{n}_{(n-2,\l)}$ is non-degenerate away from the real line (\textit{cf.} \ref{pa:gramrels}), we obtain:
\begin{corollary} Fix $l$ and let $\l\vdash l+2$. Let $C^{(\l)}\in \C[\a]$ and $P^{(\l)}_n\in \C[\a]$ be as per Theorem \ref{thm:main}, with $n\geq l+4$. Then $C^{(\l)}$ and $P^{(\l)}_n$ have real roots.   
\end{corollary}
In the cases computed thus far, we will use the mean value theorem to provide bounds on the values of roots. It will be helpful first to collect some identities for Chebyshev series with the \ramping\ property, which we do in \S\ref{ss:rampingroots}.


\subsection{Identities for \ramping\ series}\label{ss:rampingroots}
%\ppm{[could move to \S\ref{s:roots}.]}
In what follows, let $P$ be a Chebyshev series.

\mdef\label{pa:roots} 
Suppose that $P$ has the \ramping\ property for $n\geq N$ and further suppose that $P_{k}(x_0)=0$ for some $k\in \Z$ and $x_0\in \C$. Since the series has no common factors, we have that $P_{k\pm 1}(x_0)\neq 0$. The Chebyshev recursion then implies that 
\beq
    P_{k+n}(x_0)=P^U_n(x_0)P_{k+1}(x_0),  \label{eq:chebyx0}
\eq
for any $n\in \Z$. Therefore, a series with the \ramping\ property can only have a recurring (with period $N>0$) root at $x_0$, if $x_0$ satisfies $P^U_N(x_0)=0$. 

The following lemma will help us (in \S\ref{ss:roots}) analyse the distribution of the roots of the polynomials $P^{(\l)}_n$ introduced in \S\ref{ss:non1dsp}.

\begin{lemma}\label{lem:roots} 
Let $P_n$ be a series which has the \ramping\ property for $n\geq N$. Let $k, k'$, and $ r$ be positive integers with $r<k+k'$. Then
\[P_{N+k}\left(2 \cos \left(\frac{r \, \pi}{k+k'}\right)\right)= (-1)^{r} P_{N-k'}\left(2 \cos \left(\frac{r \, \pi}{k+k'}\right)\right).  \]
\end{lemma}
\begin{proof} Set $\th=\pi/(k+k')$ for this proof. Using \eqref{eq:trig} we can compute
\
\begin{align*}P^U_{k+k'}\left(2 \cos \left(r \, \th\right)\right)&= \sin\left( r \, \pi \right)/\sin\left( r \, \th \right)=0,\\
    P^U_{k+k'+1}\left(2 \cos \left(r\, \th \right)\right)&= \sin\left( r \, \pi +r \, \th \right)/\sin\left(r\,\th \right)=(-1)^r,
\end{align*}
noting that $\sin\left(r\,\th\right)\neq 0$. Therefore, by \eqref{eq:chebyx0} for any $i\in \Z $ we have
\[
P^U_{k+k'+i}\left(2 \cos \left(r\, \th\right)\right)=(-1)^r P_i^{U}\left(2 \cos (r\, \th) \right).\tag{\Celtcross}\]
By \ref{pa:chebslv}, we then have
\begin{align*}    P_{N+k}(2 \cos(r\, \th) )&=\sum_{i=-(d+1)}^{d+1} a_i P^U_{k+i}(2 \cos(r\, \th) )\\
&=(-1)^r\sum_{i=-(d+1)}^{d+1} a_i  P^U_{-k'+i}(2 \cos(r\, \th) )\\
&= (-1)^r P_{N-k'}(2 \cos(r\, \th) ),
\end{align*}
where we have used (\Celtcross) in the second line with indices shifted $i\mapsto i-k'$.
\end{proof}

\subsection{Roots of the $P^{(\l)}_n$}\label{ss:roots}

\mdef Firstly, we consider the infinite family of series $P^{(l+2)^T}_{n}(\a)$ (i.e. using the formula \eqref{eq:prow} beyond the bound for which it is verified). Parametrising by $n=l+4+k$ with $k\geq 0$, we have that $P^{(l+2)^T}_{n}(\a)$ is of degree $k+2$.  By applying Lemma \ref{lem:roots}, with $k'=2$ we obtain for any $r=1,\dots,k+1$ 
\[P^{(l+2)^T}_{n}( x^{(k)}_{r})= (-1)^r(l+2),\]
where $x^{(k)}_{r}:=2\cos(r \pi/(k+2))$. We can also determine
\begin{align*}P^{(l+2)^T}_{n}( x^{(k)}_{k+2})&=P^{(l+2)^T}_{n}( -2)=(-1)^{k+2} (4 + k + l),\\
P^{(l+2)^T}_{l+4+k}( l+2 )&=-P^{U}_{k}(l+2)<0, \quad (l>-1),\\
\lim_{x\to \infty } P^{(l+2)^T}_{n}( x)&=\lim_{x\to \infty } x^{k+2}=+\infty.
\end{align*}
In the second line, we have noted that $l+2$ is a root of $P_{l+4}^{(l+2)^T}$ and applied \eqref{eq:chebyx0}. Thus by an application of the mean value theorem we have that for $l>-1$, $P^{(l+2)^T}_{l+4+k}(\a)$ has $k+2$ distinct real roots $y^{(k)}_{r}$ with $r=0,1,\dots,k+1$ satisfying $y_0^{(k)}>l+2$ and 
\[-2<y^{(k)}_{k+1}<x^{(k)}_{k+1}<y^{(k)}_{k}<x^{(k)}_{k}< \dots <x^{(k)}_{2}<y^{(k)}_{1}<x^{(k)}_{1}<2.\]

\mdef We now consider the infinite family $P_n^{(l+2)}(\a)$ for $l>-1$. This has degree $k+3$ (with the same parametrisation $n=l+4+k$). Here we compute (for $r=1,\dots, k+1$) that
\begin{align*}
    P_n^{(l+2)}(\a)(x_r^{(k)})&=(-1)^r(l+2)(\a+2 l),\\
    P_n^{(l+2)}(\a)(x_0^{(k)})&=2 (1 + l) (2 + l). 
\end{align*}
For $l>0$, the factor $(\a+2 l)$ is positive on $(-2,2)$. Therefore, in this case, the mean value theorem gives that $P_{l+4+k}^{(l+2)}(\a)$ has at least $k+1$ distinct roots which may be written as $z^{(k)}_{r}$ with $r=0,1,\dots,k,$ satisfying
\[-2<x_{k+1}^{(k)}<z_k^{(k)}<x_{k}^{(k)}< \dots<x_1^{(k)}<z^{(k)}_{1}<x_1^{(k)}<z^{(k)}_{0}<x_0^{(k)}=2.\]
The distribution of remaining roots (if they exist) depends on the values of $k$ and $l$. For instance, when $l>1$ and $n=l+4+k>6$ (all but finitely many cases when $l>1$), $P_{l+4+k}^{(l+2)}(\a)$ has two further roots $z^{(k)}_{k+1}$ and $z^{(k)}_{k+2}$ which satisfy
\[-\infty< z^{(k)}_{k+2}<- 2 (l+1) <- 2 l <z^{(k)}_{k+1}<-(l+1).\]
These bounds can be determined by evaluating $P_{l+4+k}^{(l+2)}$ at the (integral) roots of the $k=-1,-2,$ cases (\textit{c.f} remark \ref{rem:polys}) using \eqref{eq:chebyx0}, and by considering the behaviour as $|x|\to \infty$.

\mdef We can use similar arguments to obtain bounds on roots for the series $P^{(l+1,1)}_{n}$ and $P^{(l+1,1)^T}_{n}$ also. Here $P^{(l+1,1)}_{l+4+k}$ has degree $k+5$ and $P^{(l+1,1)^T}_{l+4+k}$ has degree $k+4$. Noting that $P^{(l'+1,1)}_{l'+3}(\a)$ is negative on $(-2,2)$ when $l'>4$, and $P^{(l+1,1)}_{l+3}(\a)$ is positive on $(-2,2)$ when $l>3$, we apply Lemma \ref{lem:roots} (now with $k'=1$) to obtain 
\[
    \text{sgn}(P^{(l'+1,1)}_{l'+4+k}(X_r^{(k)}))=-(-1)^r=
    -\text{sgn}(P^{(l+1,1)^T}_{l+4+k}(X_r^{(k)})),
\]  
where $X_r^{(k)}:=2\cos(r \pi/(k+1))$ with $r=1,\dots,k$, for $l'$ and $l$ as such. Thus in each case we observe $k-1$ distinct roots, $Y_r^{(k)}\in (-2,2)$ with $r=1,\dots,k-1$, satisfying
\[-2<X_{k}^{(k)}<Y_{k-1}^{(k)}<X_{k-1}^{(k)}< \dots<X_2^{(k)}<Y^{(k)}_{1}<X_1^{(k)}<2.\]
We observe another two roots $X_1^{(k)}<Y_0^{(k)}<X_{0}^{(k)}=2$ and $-2=X_{k+1}^{(k)}<Y_{k}^{(k)}<X_{k}^{(k)}$ in each case by explicitly verifying
\begin{align*}
    -(-1)^0P^{(l'+1,1)}_{l'+4+k}(X_0^{(k)})&>0, &
    -(-1)^{k+1} P^{(l'+1,1)}_{l'+4+k}(X_{k+1}^{(k)})&>0,\\
    (-1)^{0} P^{(l+1,1)^T}_{l+4+k}(X_{0}^{(k)})&>0, &
    (-1)^{k+1} P^{(l+1,1)^T}_{l+4+k}(X_{k+1}^{(k)})&>0,
\end{align*}
for $l,l'$ as before and $k\geq 0$. 

This leaves (at most) four roots of $P^{(l'+1,1)}_{n}$ unaccounted for, and three of $P^{(l+1,1)^T}_{n}$. These remaining roots can be bounded by evaluating at (integral) roots of the $k=-1,-2$ cases (\textit{c.f} Remark \ref{rem:polys}), and by considering the asymptotic behaviour. We find the following: $P^{(l+1,1)^T}_{n}$ (with $l>3$) has three further roots $Y_{k+1}^{(k)},Y_{-1}^{(k)},Y_{-2}^{(k)}$ which satisfy
\[Y_{k+1}^{(k)}<-2<2<l-1<Y_{-1}^{(k)}<l<l+1<Y_{-2}^{(k)},\]
and $P^{(l'+1,1)}_{n}$ (with $l'>4$) has four further roots $Z_{k+3}^{(k)},Z_{k+2}^{(k)},Z_{k+1}^{(k)},Z_{-1}^{(k)}$ which satisfy
\[Z_{k+3}^{(k)}<-2l'< Z_{k+2}^{(k)}<-l'+1<-l'+2<Z_{k+1}^{(k)}<-l'+3\leq -2<2<Z_{-1}^{(k)}.\]
The roots of the series $P^{(2,2)}_{6+k}(\a)$ can be analysed similarly, but we omit this for brevity. Based on our observations here, we conjecture:

\begin{conjecture} Fix $l$ and let $\l \vdash l+2$. For each $k\geq -1$, the polynomial $P^{(\l)}_{l+4+k}$ is square free. Furthermore, for large enough $k$, $P^{(\l)}_{l+4+k}$ has a fixed number of roots lying outside the interval $[-2,2]$, and the set
\[\bigcup_{k=0}^{\infty } \{\a \in [-2,2] \ | \ P^{(\l)}_{l+4+k}(\a)=0  \},\]
    is dense \ppmm{in $[-2,2]$}.
\end{conjecture}


\section{Gram-determinant-decorated Rollet Graphs}   \label{ss:gdetdecrollet}

Theorem~\ref{thm:main} is the main technical result in the computation of the `bootstrap' data 
for the ToR method for representation theory of KY algebras. 
In \S\ref{ss:repthy} we will see how to apply this directly in representation theory. 
But it is also interesting from a general representation theory point of view to consider KY gram matrices beyond the bootstrap. In the $l=-1$ case there is a very neat pattern, related to the corresponding classical problem for symmetric groups over fields of finite characteristic. 
Here we will evidence that this pattern generalises to other $l$ values. 

In \S\ref{ss:drgl-1} and \S\ref{ss:gdTL} we review the $l=-1$ case. 
From this we see that it is natural to express $\Delttt{n}{p,\l}$ in the form of the corresponding marginal vertex function $\V^n_{(p,\l)}$; and that this `correction' factor is related to ratios of the appropriate Chebyshev polynomials raised to a suitable power. 
In \S\ref{ss:e-next!} we  write this in a way that 
formally generalises to all $l$. 
In \S\ref{ss:gdetl0}-\ref{ss:gdetl2} we compare this to the results of direct computation. 



\newcommand{\ournotmain}{arm}

\subsection{Statement of \ournotmain\ property and conjecture \ref{conj:onm0}}
\label{ss:e-next!}

%\ppm{[Delete the following list now?:]}\bajm{[done]}

\newcommand{\CC}{{\mathcal C}}  %% hook for claimed marginal factors

\mdef   \label{de:arm-property}
%\bajm{[to quote radiohead: ``I dont belong here"]} \ppm{[I must say, I am pretty lost here myself at the moment. Probably we need a quick chat about what should be in this Section and what should be later?]}
%\bajm{[I will check this to make sure its right]}
Armed with the polynomials $P_n^{(\l)}$, determined in \S\ref{ss:non1dsp}, 
\ppmm{and motivated by the proof of the $l=-1$ formula \eqref{eq:det-recurs} (and finite rank calculations),}
we can define 
\beq  \label{eq:defCC} 
\CC^{\ppmm{n}}_{(p,\lambda)} 
= \prod_{{(p+1,\mu )\; adjacent \; \atop to \; (p,\lambda) \ in \ \Roll_l}} 
     \left( \frac{P_n^{(\mu)}}{P_{n-1}^{(\mu)}} \right)^{dim (\Deltaa^{n-1}_{(p+1,\mu )}) } 
\eq 
- at least for values of $(p,\lambda)$ where the graphical condition makes sense. 
%\ppm{ [Not finished - need to lock n to p etc. And deal with $\lambda$s not in a branch. Idea is that in the end $\V^n_{(p,\lambda)}$ is given by $\CC^{n}_{(p,\lambda)}$.]}

The utility of $\CC^n_{(p,\l)}$ comes from the range of values of the index for which we have 
an identity $\V^n_{(p,\lambda)} = \CC^n_{(p,\lambda)}$ (cf. \ref{eq:MVF1}). 
We call this {\em the \ournotmain\ property}. 
See in particular Conjecture \ref{conj:onm0}. 

\medskip 

A key theorem (extended by a conjecture) we consider in this paper relates the marginal vertex function 
\eqref{eq:MVF1}, 
on the arms of our Rollet graphs 
\ref{de:Rollet}, 
to a family of Chebyshev polynomials defined for each integer partition. 
In other words it gives a way, in principle, to compute almost all KY gram determinants. 

%\ppm{(read the following [...] before deleting!:)[Again I would not do it in this order. Minimize the amount of conjectures. Let's try to state all we can prove as Proposition - writing it as a general formulation and then saying it holds in the given cases. Then at the end conjecture that the general formulation works generally.]}

  \begin{theorem}[\bajmm{The arm property}]\label{thm:V=C}
      %/Conjecture}}. 
%Fix  $-1\leq l\leq 1$.   %%%% Being brace here!!...
Sufficient conditions on $l,m,$ and $p$, such that 
for each  $\l \vdash l+2$, 
%there exists \ppm{[let's say what they are?]}
%a family of monic, integral, polynomials $P^{(\l)}_k$ for $k\geq l+2$ satisfying the Chebyshev recursion {as in (\ref{eq:Cheby0})}   such that, 
%
%the gram determinants of 
the arm $J_{l,p+2m}$ standard modules $\Deltaaa{p+2m}{(p,\l)}{}$ with  $p\geq l+2$ and $m\geq 1$, 
\ppmm{obey the \ournotmain\ property --- that is,   
the gram determinants,} 
%are given by  
satisfy 
%\bajm{[satisfy seems more appropriate to me]}
\beq   \label{eq:V=C} 
\V^{n=p+2m}_{(p,\l)} = \CC^{n=p+2m}_{(p,\l)}
\eq 
as per (\ref{eq:MVF1}) and (\ref{eq:defCC})
--- are as follows. For $m=1$ the identity holds for all $l$ and $p\geq l+3$.
For $l=-1$ the identity holds for all $n=p+2m$. 
For $l=0$ the identity holds when $m=2$ and $p\leq 8$, and when $m=3$ and $p\leq 5$. For $l=1$ the identity when $m=2$ and $p\leq 7$.
%\ppm{[i made these bounds up when drafting. what are the correct ones?]}\bajm{[I'll add them]}\ppm{[perhaps need p/m conditions as well?]}
%That is 
Since we are in the arm, 
equation (\ref{eq:V=C}) is the same as 
to say
\[  %\frac{
\Delttt{(n)}{p,\l}   
\;\; = \; \left(\frac{P^{(\l)}_{n}}{P^{(\l)}_{n-1}}\right)^{d^{{n-1}}_{(p+1,\l)}} %}{
\prod_{v \text{ adjacent}\atop \text{to } (p,\l) \text{ in } \Roll_{l}} \!\! \Delttt{(n-1)}{v}    
%} = 
\;\; ,
\]
%\ppm{[-so far I don't see this directly agreeing with above! - just something for me to check. I guess the point is the $n\geq l+2$ assumption, but so far this doesn't seem like the right way of putting it to me.]}\bajm{[Actually the $n\geq l+2$ assumption is no assumption at all, the $p\geq l+2$ is the point, this ensures we are on the tails of the rollet which is the only place the conjecture we originally had holds]}
where 
$$
%d^{{n-1}}_{(p+1,\l)}  \; :=  \; \text{dim}(\Deltaaa{n-1}{(p+1,\l)})    . 
d^{{n}}_{(p,\l)}  \; :=  \; \text{dim}(\Deltaaa{n}{(p,\l)})    . 
$$
  \end{theorem}

\begin{proof} For $m=1$, and $p\geq l+3$, \eqref{eq:V=C} becomes
\[\frac{\Delttt{(p+2)}{p,\l}}{\Delttt{(p+1)}{p-1,\l}}=\left(\frac{P^{(\l)}_{n}}{P^{(\l)}_{n-1}}\right)^{d^{p+1}_{(p+1,\l)}}, \]
which follows from Theorem \ref{thm:main} since $d^{p+1}_{(p+1,\l)}=d_\l$.

For the case $l=-1$ note from \eqref{eq:det-recurs} that we have 
the \ournotmain\ property on the nose for all $n$. For the other cases the identity holds by direct calculation over the indicated ranges. 
\end{proof}


\begin{conjecture}   \label{conj:onm0} Formula \eqref{eq:V=C} holds for all $l\geq -1$, $p\geq l+2$ and $m\geq 1$.
\end{conjecture}

%\bajm{[Delete the following conjecture?]}
\ignore{\begin{conjecture}   \label{conj:onm01}
Fix integers $l\geq -1$ and $n \geq l+2$. For each integer partition $\l \vdash l+2$, there exists 
\ppm{[-well we have already had a stab at saying what they are. if we retract that, this is non-empty but weaker. And the $\CC$ formulation would need adjustment.]}
a family of monic, integral, polynomials $P^{(\l)}_k$ for $k\geq l+2$ satisfying the Chebyshev recursion {as in (\ref{eq:Cheby0})}  %\bajm{\sout{[refer to something...?]}} 
such that, the gram determinants of the $J_{l,n}$ standard modules, $\Delta^{(n)}_{(p,\l)}$ with  $p\geq l+2$, satisfy 
\[
\V^{n}_{(p,\l)} = \CC^{n}_{(p,\l)}
\]
as in (\ref{eq:MVF1}) and (\ref{eq:defCC}), 
that is to say
\[
\frac{|\Delta^{(n)}_{(p,\l)}|}{\prod_{v \text{ adjacent}\atop \text{to } (p,\l) \text{ in } \Roll_{l}} |\Delta^{(n-1)}_{v}|}
= \left(\frac{P^{(\l)}_{p+1}}{P^{(\l)}_{p}}\right)^{d^{\ppm{n-1}}_{(p+1,\l)}},
\]
\ppm{[-so far I don't see this directly agreeing with above! - just something for me to check. I guess the point is the $n\geq l+2$ assumption, but so far this doesn't seem like the right way of putting it to me.]}\bajm{[Actually the $n\geq l+2$ assumption is no assumption at all, the $p\geq l+2$ is the point, this ensures we are on the tails of the rollet which is the only place the conjecture we originally had holds]}
where 
$d^{\ppm{n-1}}_{(p+1,\l)}=\text{dim}(\Delta^{(n-1)}_{(p+1,\l)})$.
\ppm{[-RHS should depend on same things as LHS (unless really indep.).]}

\end{conjecture}}


\subsection{On proof of proven part}   \label{ss:gdetdecrolletl-1}
%\bajm{[section name?] [Delete this subsection?]}
%...


%\subsection{Examples} 


\subsubsection{Decorated Rollet Graph, for $l=-1$}  \label{ss:drgl-1}

See Fig.\ref{fig:gramdets0312} for the Rollet graph $\Roll_{-1}$ augmented with its fibres of standard gram determinants.
And 
(for the impatient reader, jump ahead to)
see Fig.\ref{fig:gramdets03} for a schematic for the representation theory interpretation. 

\begin{figure} 
%
\newcommand{\du}[2]{\begin{array}{c} \textcolor{red}{{#1}} \\ \hspace{-.321in}{#2}\hspace{-.321in} \end{array}}
\newcommand{\sqrr}{\square\ar[r]}
%
%
\[ { } \hspace{-.253in}
%
\xymatrix @C=15pt @R=1pt {
  {\emptyset}\ar[r] & \sqrr & \sqrr & \sqrr & \sqrr & \sqrr & \sqrr & \sqrr & \sqrr & \sqrr & { }
  \\
  {\lambda=0  }       &  1             &  2      &  3  &  4  &  5   & 6 & 7 & 8 & 9 
  \\  %\hline
    {\du{1}{1} }\ar[dr]
    \\
& {\du{1}{1}}\ar[dr]\ar[dl]^{2/1}
\\ 
{\du{1}{{[2]}} }\ar[dr]    &                  &
{\du{1}{1}}\ar[dr]\ar[dl]^{3/2}
\\
&{\du{2}{[3]}} \ar[dr]\ar[dl]^{2/1}               &     &
{\du{1}{1}}\ar[dr]\ar[dl]^{4/3}
\\ 
{ \du{2}{{[ 2]}^2 [3]}}  \ar[dr]  &                     &
{\du{3}{[4]} }  \ar[dr]\ar[dl]^{3/2}  &    &
{\du{1}{1} } \ar[dr]\ar[dl]^{5/4}
\\
    & {\du{5}{{\frac{ [3]^4 [4]}{[2]}^{}}}}  \ar[dr]\ar[dl]^{2/1}   &  &
    {\du{4}{[5]}} \ar[dr]\ar[dl]^{4/3} &  &
    {\du{1}{1}} \ar[dr]\ar[dl]^{6/5}
    \\ 
    {\du{5}{{[2]}^4 [3]^4 [4]}} &              &
    {\du{9}{{\frac{[4]^5[5]}{[2]}^{}}}} &    &
    {\du{5}{[6]}} & &{ \du{1}{1} }
    \\ 
    &{\du{14}{\frac{[3]^{13}[4]^6[5]}{[2]^{6}}}} &&
    {\du{14}{\frac{[5]^6[6]}{[2]^{}}}} &&
    {\du{6}{[7]}} &&
    {\du{1}{1}}
 \\
   {\du{14}{[2]^8[3]^{13}[4]^6[5]}}  &&
   {\du{28}{\frac{[4]^{20}[5]^7[6]}{[2]^{7}[3]^{}}}} &&
   {\du{20}{\frac{[6]^7[7]}{[2]^{}}}} &&
   {\du{7}{[8]}}  &&
   {\du{1}{1}} \\
   & {\du{42}{\frac{[3]^{40}[4]^{26}[5]^8[6]}{[2]^{27}}}} &&
   {\du{48}{\frac{[5]^{27}[6]^8[7]}{[2]^{8}[3]^{}}}} &&
   {\du{27}{\frac{[7]^8[8]}{[2]^{}}}} &&
   {\du{8}{[9]}} &&
   {\du{1}{1}}
 \\
   {\du{42}{[2]^{15}[3]^{40}[4]^{26}[5]^{8}[6]}}\ar[dr] &&
   {\du{90}{\frac{[4]^{74}[5]^{35}[6]^{9}[7]}{[2]^{35}[3]^{9}}}}\ar[dl]^{3/2} &&
   {\du{75}{
   %[4]^{0}[5]^{0}
       \frac{  [6]^{35}[7]^{9}[8]}{[2]^{9}[3]^{}}}} &&
   {\du{35}{\frac{[8]^9[9]}{[2]^{}}}} &&
   {\du{9}{[10]}} &&
   {\du{1}{1}}
   \\
   & { \begin{array}{c} { } \\ \textcolor{red}{132} \\ 
        {[2]^{15-90-35}} \\ {[3]^{40+90-9}} \\
        {[4]^{26+74}} \\ {[5]^{8+35}} \\ {[6]^{1+9}[7]}
        \end{array} }
}
\]


\caption{The $l=-1$ Rollet graph (top) augmented with many standard gram determinants (parametrised as in Fig \ref{fig:gramdets03}) --- organised with increasing rank $n$ going downwards, each under their fixed standard module label. Edges amongst these fibres of nodes indicate schematically how they are computed from \cite{Martin91}. \label{fig:gramdets0312}}
\end{figure}

\ignore{{
\begin{figure} 
%
\newcommand{\du}[2]{\begin{array}{c} \textcolor{red}{{#1}} \\ \hspace{-.321in}{#2}\hspace{-.321in} \end{array}}
\newcommand{\sqrr}{\square\ar[r]}
\newcommand{\goll}{\ar@/_1.3pc/[ll]} 
\newcommand{\gollll}{\ar@/_1.4pc/[llll]} 
\newcommand{\gollllll}{\ar@/_1.5pc/[llllll]} 
\newcommand{\gollllllll}{\ar@/_1.7pc/[llllllll]} 
%
%
\[ { } \hspace{-.73in} %\hspace{-.253in}
%
\xymatrix @C=15pt @R=1pt {
  {\emptyset}\ar[r] & \sqrr & \sqrr & \sqrr & \sqrr & \sqrr & \sqrr & \sqrr & \sqrr & { }
  \\
  {\lambda=0  }       &  1             &  2      &  3  &  4  &  5   & 6 & 7 & 8 
  \\  %\hline
    {\du{1}{1} }\ar[dr]
    \\
& {\du{1}{1}}\ar[dr]\ar[dl]^{2/1}
\\ 
{\du{1}{{[2]}} }\ar[dr]    &                  &
{\du{1}{1}}\ar[dr]\ar[dl]^{3/2}
\\
&{\du{2}{[3]}} \ar[dr]\ar[dl]^{2/1}               &     &
{\du{1}{1}}\ar[dr]\ar[dl]^{4/3}
\\ 
{ \du{2}{{[ 2]}^2 [3]}}  &                     &
{\du{3}{[4]} }&    &
{\du{1}{1} }
\\
    & {\du{5}{{\frac{ [3]^4 [4]}{[2]}^{}}}} &  &
    {\du{4}{[5]}} &  &
    {\du{1}{1}}\ar@/_1.3pc/[ll]
    \\ 
    {\du{5}{{[2]}^4 [3]^4 [4]}} &              &
    {\du{9}{{\frac{[4]^5[5]}{[2]}^{}}}} &    &
    {\du{5}{[6]}} & &{ \du{1}{1}}\ar@/_1.4pc/[llll]
    \\ 
    &{\du{14}{\frac{[3]^{13}[4]^6[5]}{[2]^{6}}}} &&
    {\du{14}{\frac{[5]^6[6]}{[2]^{}}}} &&
    {\du{6}{[7]}}\ar@/_1.3pc/[ll] &&
    {\du{1}{1}}\gollllll
 \\
   {\du{14}{[2]^8[3]^{13}[4]^6[5]}}  &&
   {\du{28}{\frac{[4]^{20}[5]^7[6]}{[2]^{7}[3]^{}}}} &&
   {\du{20}{\frac{[6]^7[7]}{[2]^{}}}} &&
   {\du{7}{[8]}}\gollll  &&
   {\du{1}{1}}\gollllllll \\
   & {\du{42}{\frac{[3]^{40}[4]^{26}[5]^8[6]}{[2]^{27}}}} &&
   {\du{48}{\frac{[5]^{27}[6]^8[7]}{[2]^{8}[3]^{}}}} &&
   {\du{27}{\frac{[7]^8[8]}{[2]^{}}}}\goll &&
   {\du{8}{[9]}}\gollllll &&
   {\du{1}{1}}
 \\
   {\du{42}{[2]^{15}[3]^{40}[4]^{26}[5]^{8}[6]}} &&
   {\du{90}{\frac{[4]^{74}[5]^{35}[6]^{9}[7]}{[2]^{35}[3]^{9}}}} &&
   {\du{75}{
   %[4]^{0}[5]^{0}
       \frac{  [6]^{35}[7]^{9}[8]}{[2]^{9}[3]^{}}}} &&
   {\du{35}{\frac{[8]^9[9]}{[2]^{}}}}\gollll &&
   {\du{9}{[10]}}\gollllllll &&
   {\du{1}{1}}\goll
}
\]


\caption{The $l=-1$ Rollet graph augmented with many standard gram determinants --- 
and curved arrows indicating standard module morphisms in case $[5]=0$. 
\label{fig:gramdets03}}
\end{figure}
}}

Figure~\ref{fig:gramdets0312} is already quite dense. Let us use this case to illustrate how the information can be presented more concisely using the marginal vertex function. 




\mdef \label{th:not-main--1}
For neatness set $d_i = \dim(\Delta_{(i,(1))}^{(n-1)})$. 
Then  $\RollV_{-1}$ is given by  
\[
\begin{tikzcd}[row sep=tiny,column sep=tiny]
    \a^{d_1}\ar[r,dash]&\left(\frac{(\alpha -1) (\alpha +1)}{\alpha }\right)^{d_2} \ar[r,dash] &
    \left(\frac{\alpha  \left(\alpha ^2-2\right)}{(\alpha -1) (\alpha +1)}\right)^{d_3}\ar[r,dash] &
    \left(\frac{\left(\alpha ^2-\alpha -1\right) \left(\alpha ^2+\alpha -1\right)}{\alpha  \left(\alpha^2 -2\right)}\right)^{d_4} \ar[r,dash] &
    \left(\frac{  %(\alpha -1) \alpha  (\alpha +1) 
    \alpha  (\alpha^2 -1) 
    \left(\alpha ^2-3\right)}{\left(\alpha ^2-\alpha -1\right) \left(\alpha ^2+\alpha -1\right)}\right)^{d_5} \ar[r,dash] &
    \dots
\end{tikzcd}
\]
- the idea is that this captures each fibre complete:
each $d_i$ requires an $n$ as input, so the expression gives the fibre as $n$ varies. 
That is, to recover the fibre innards one plugs in $n$ and $p$ to work out $m$.  

A proof can be found in, for example \cite[Ch.10]{paulsonlinelecturenotes}, 
and we include the relevant extract here in \S\ref{ss:gdTL}.  
In this setting (classical) Chebyshev polynomials arise as part of the geometry of Young diagrams (hook lengths and so on - the machinery of the quantisation of Young's normal forms), or equivalently of walks on the Bratelli diagram. 

%(although 
(While there are certainly other possible methods, and it is far from clear if the strategy there offers the best chance of a lifting to other $l$,
this method is certainly suggestive of a generalisation, as we will see.) 


\mdef  
For example the fixed $m = (n-i)/2$ slice with $m=1$ is: 
\ignore{{
Let us decorate a rollet graph 
\ppm{[it feels like more words are needed here!? for the moment I am confused by what is here. doesn't seem quite right to me. Is this intended to be $\RollV_{-1}$, or $\Roll^{(m)}_{-1}$ for some $m$, or something else?]}
as per the $l=-1$ case
}}
\[
\begin{tikzcd}[row sep=tiny,column sep=tiny]
    \a^2\ar[r,dash]&\frac{(\alpha -1)^3 (\alpha +1)^3}{\alpha ^3}\ar[r,dash]&\frac{\alpha ^4 \left(\alpha ^2-2\right)^4}{(\alpha -1)^4 (\alpha +1)^4}\ar[r,dash]&\frac{\left(\alpha ^2-\alpha -1\right)^5 \left(\alpha ^2+\alpha -1\right)^5}{\alpha ^5 \left(\alpha ^2-2\right)^5}\ar[r,dash]&\frac{(\alpha -1)^6 \alpha ^6 (\alpha +1)^6 \left(\alpha ^2-3\right)^6}{\left(\alpha ^2-\alpha -1\right)^6 \left(\alpha ^2+\alpha -1\right)^6}\ar[r,dash]&\dots
\end{tikzcd}
\]


%\pagebreak
\subsection{Examples} 

\subsubsection{Decorated Rollet Graphs for $l=0$}   \label{ss:gdetl0}

Here we provide some examples of decorated Rollet graphs for $l=0$ to verify Theorem \ref{thm:V=C}. Firstly, recall the undecorated Rollet graph for $l=0$ is as follows:
\[
\begin{tikzcd}[row sep=tiny,column sep=tiny]&&(2)\ar[r,dash]&(2)\ar[r,dash]&(2)\ar[r,dash]&\dots\\
    \emptyset\ar[r,dash]&(1)\ar[ur,dash]\ar[dr,dash]\\
    &&(1^2)\ar[r,dash]&(1^2)\ar[r,dash]&(1^2)\ar[r,dash]&\dots
\end{tikzcd}
\]

\mdef   \label{pa:l0OEIS}
We will be using the standard module dimensions again. 
The table of standard module dimensions for $l=0$ starts as in Fig.\ref{fig:l=0Brat}. 
%\bajm{[delete the follow sentence?]}For example, the sequence of spine dimensions starts:
%1, 1, 3, 11, 43, 173, 707, ....
%This, apart from the leading 1, is A026671 in \cite{OEIS}. 
%We will return to this point \ppm{ in \ref{where???!}. }


\begin{figure} 
\[
\xymatrix@R=-2pt@C=7pt{
1 \ar@[red][ddddrr] \\ \mbox{ } \vspace{.1in} \; \\ \\ \\ 
&& 1 \ar[dddll] \ar[ddrr] \ar[ddddr] \\ \\ 
 &&&& 1 \ar@[red][ddddddll] \ar@[red][dddddrr] \\
1 \ar@[red][dddddrr] \\
 &&& 1 \ar@[red][ddddl] \ar@[red][dddddrr] \\ \\ \\ 
 &&&&&&1 \ar[ddddrr] \ar[ddddll] \\
 &&3 \ar[ddddll] \ar[dddrr] \ar[dddddr] \\ 
 &&&&&1 \ar[ddddrr] \ar[ddddll] \\ \\ 
 &&&&4 \ar@[red][dddddll] \ar@[red][ddddrr] &&&& 1 \ar@[red][ddddrr] \ar@[red][ddddll] \\
 3 \ar@[red][ddddrr] \\
 &&& 4 \ar@[red][dddl] \ar@[red][ddddrr] &&&& 1 \ar@[red][ddddrr] \ar@[red][ddddll] \\ 
 \\ 
 &&&&&&5 \ar[ddddrr] \ar[ddddll] &&&& 1 \ar[ddddrr] \ar[ddddll] \\
&& 11 \ar[ddddll] \ar[dddrr] \ar[dddddr] \\  
&&&&& 5 \ar[ddddrr] \ar[ddddll] &&&& 1 \ar[ddddrr] \ar[ddddll] \\ 
\\ 
&&&& 16 \ar@[red][dddddll] \ar@[red][ddddrr] &&&& 6 \ar@[red][ddddrr] \ar@[red][ddddll] &&&& 1 \ar@[red][ddddrr] \ar@[red][ddddll] \\
11 \ar@[red][ddddrr] \\
&&& 16 \ar@[red][dddl] \ar@[red][ddddrr] &&&& 6 \ar@[red][ddddrr] \ar@[red][ddddll] &&&& 1 \ar@[red][ddddrr] \ar@[red][ddddll] \\ 
\\ 
 &&&&&& 22 \ar[ddddrr] \ar[ddddll] &&&& 7 \ar[ddddrr] \ar[ddddll] &&&& 1 \ar[ddddll] \ar[ddddrr] \\
&& 43 \ar[ddddll] \ar[dddrr] \ar[dddddr] \\  
&&&&&  22 \ar[ddddrr] \ar[ddddll] &&&& 7 \ar[ddddrr] \ar[ddddll] &&&& 1 \ar[ddddll] \ar[ddddrr] \\ 
\\ 
&&&& 65 &&&& 29 &&&& 8 &&&& 1 \\
43 \\
&&& 65 &&&& 29 &&&& 8 &&&& 1 \\
}
\]
\caption{ Full Bratelli diagram for $l=0$. Edges from an even $n$ to an odd $n$ layer are coloured red, just to aid visibility of layers. \label{fig:l=0Brat} }
%\bajm{[is it worth adding dashed horizontal lines to separate the different $n$-values?]} \ppm{[yes. that's what the colour is for, but it is not super clear. it might also be better to use the LR symmetric layout? - which shows the relevance of the OEIS sequence!!]}}
\end{figure} 


\begin{figure}
%\newcommand{\baa}[1]{\!\!\begin{array}{c} #1 \end{array}\!\!}

\vspace{-4.4cm} 
\[ \hspace{-4.7cm} \vspace{-1cm} 
\xymatrix@R=1pt@C=-20pt{
\emptyset && (1) \ar[ll]^x & (2) \ar[l]^{\frac{(x-1)(x+2)}{2x}} 
& (1^2) \ar@/_1.3pc/[ll]_{\frac{x-1}{2}} 
& (2)  \ar@/^1.3pc/[ll]^{\frac{x(x^2+x-4)}{(x-1)(x+2)}}
& (1^2) \ar@/_1.3pc/[ll]_{\frac{(x+1)(x-2)}{x-1}}   \\ 
\\  
1 \ar@[red][ddddddrr] \\  \mbox{ } \vspace{.41in} \; \\ \\ \\  \\ \\ 
&& 1 \ar[dddll]^{x} \ar[ddrr] \ar[ddddr] \\ \\ 
 &&&& 1 \ar@[red][ddddddll] \ar@[red][dddddrr] \\
{\baa{x \\ \ppm{1}}} \ar@[red][dddddrr] \\
 &&& 1 \ar@[red][ddddl] \ar@[red][dddddrr] \\ \\ \\ 
 &&&&&&1 \ar[ddddrr] \ar[ddddll] \\
 && {\baa{(x-1)^2 \\  (x+2) \\ \ppm{3}}} \ar[ddddll]^{x} \ar[dddrr] \ar[dddddr] \\ 
 &&&&&1 \ar[ddddrr] \ar[ddddll] \\ \\ 
 &&&&{\baa{(x-1)(x+1) \\ (x-2)(x+2) \\ \ppm{4}}} \ar@[red][dddddll] \ar@[red][ddddrr] &&&& 1 \ar@[red][ddddrr] \ar@[red][ddddll] \\
 {\baa{x^3(x-1)^2 \\ (x+2) \\ \ppm{3}}} \ar@[red][ddddrr] \\
 &&& {\baa{x(x-1) \\ (x^2+x-4) \\ \ppm{4}}} \ar@[red][dddl] \ar@[red][ddddrr] &&&& 1 \ar@[red][ddddrr] \ar@[red][ddddll] \\ 
 \\ 
 &&&&&&{\baa{(x-1)(x+2) \\ (x^3-\!x^2-\!3x+\!1) \\ \ppm{5}}} \ar[ddddrr] \ar[ddddll] &&&& {\;\;1\;\;} \ar[ddddrr] \ar[ddddll] \\
&& {\baa{(x-1)^{12}(x+1) \\ (x-2)  (x+2)^6 \\ (x^2+x-4) \\ \ppm{11}}} \ar[ddddll]^x \ar[dddrr] \ar[dddddr] \\  
&&&&& {\baa{(x-1)\\ (x^4\!+\!x^3\!-\!5x^2\!-\!x\!+\!2) \\ \ppm{5}}} \ar[ddddrr] \ar[ddddll] &&&& {\;\;\;\;\;\;1\;\;\;\;\;\;} \ar[ddddrr] \ar[ddddll] \\ 
\\ 
&&&& {\baa{(x\!-\!1)^8(x\!+\!1)^6(x\!-\!2)^6 \\ (x+2)^7(x^2+x-4)^{} \\ (x^3-x^2-3x+1) \\ \ppm{16}}} \ar@[red][dddddll] \ar@[red][ddddrr] &&&& {\baa{(x-\!1)^2 (x\!+\!2) \\ (x^3-4x-2) \\ \ppm{6}}} \ar@[red][ddddrr] \ar@[red][ddddll] &&&& {\;\;\;\;1\;\;\;\;} \ar@[red][ddddrr] \ar@[red][ddddll] \\
{\baa{x^{11}(x-1)^{12} \\ (x+1) (x-2) \\ (x+2)^6 (x^2+\!x-\!4) \\ \ppm{11}}} \ar@[red][ddddrr] \\
&&& {\baa{x^5(x\!-\!1)^8(x\!+\!\!1)(x\!-\!2) \\ (x+2) (x^2+x-4)^{\bajm{6} }\\ (x^4\!+\!x^3\!-\!5x^2-\!x+\!2) \\ \ppm{16}}} \ar@[red][dddl] \ar@[red][ddddrr] &&&& {\baa{(x-1)^2 x \\ (x^3\!+\!2x^2\!-\!4x\!-\!6) \\ \ppm{6}}} \ar@[red][ddddrr] \ar@[red][ddddll] &&&& {\;\;\;\;1\;\;\;\;} \ar@[red][ddddrr] \ar@[red][ddddll] \\ 
\\ 
 &&&&&& \ppm{22} \ar[ddddrr] \ar[ddddll] &&&& \ppm{7} \ar[ddddrr] \ar[ddddll] &&&& {\;\;1\;\;} \ar[ddddll] \ar[ddddrr] \\
&& \ppm{43} \ar[ddddll] \ar[dddrr] \ar[dddddr] \\  
&&&&&  \ppm{22} \ar[ddddrr] \ar[ddddll] &&&& \ppm{7} \ar[ddddrr] \ar[ddddll] &&&& {\;\;1\;\;} \ar[ddddll] \ar[ddddrr] \\ 
\\ 
&&&& \ppm{65} &&&& \ppm{29} &&&& \ppm{8} &&&& {\;\;1\;\;} \\
\ppm{43} \\
&&& \ppm{65} &&&& \ppm{29} &&&& \ppm{8} &&&& {\;\;1\;\;} \\
}
\]
%
    \caption{Full Bratelli for $l=0$ giving gram determinants (with module dimensions below in red). 
    The top row indicates the conjectured ($n$-independent up to dimensions) edge factors. }
    \label{fig:Bratelli-l0-gramdets}
\end{figure}




\mdef 
The full Bratelli diagram for $l=0$ giving gram determinants (with dimensions below in red) 
starts as in Fig.\ref{fig:Bratelli-l0-gramdets}. 
    The very top row in the figure indicates the conjectured ($n$-independent up to dimensions giving exponents) edge factors. 

\begin{proposition}\label{pr:1cupgram}
The 1-cup gram-det-decorated graph $\Roll^1_0$ (as in \ref{de:Roll-etc}) is   %\ppm{[fix \;it-]}
\[  \hspace{-1.5cm} 
\begin{tikzcd}[row sep=tiny,column sep=tiny]
&&  \overbrace{\alpha (\alpha-1)(\alpha^2 +x-4)}^{\Delttt{4}{2,(2)} \;=\; C^{(2)}P^{(2)}_4} \ar[r,dash] &
  \overbrace{(\alpha-1)(\alpha^4 +\alpha^3 -5x^2 -x+2)}^{C^{(2)}P^{(2)}_{5}} \ar[r,dash] &
  C^{(2)}P^{(2)}_{n=6} \ar[r,dash] & \dots
  \\
    \alpha \ar[r,dash] &
    \overbrace{
    (\alpha-1)^2 (\alpha+2)   }^{\Delttt{3}{1,(1)}  }   \ar[ur,dash]\ar[dr,dash]
    \\
    &&\underbrace{(\alpha^2 -1)(\alpha^2 -4)}_{C^{(1^2)}P^{(1^2)}_4} \ar[r,dash] &
    \underbrace{(\alpha-1)(\alpha+2)(\alpha^3 -\alpha^2 -3\alpha+1)}_{C^{(1^2)}P^{(1^2)}_5} \ar[r,dash] &
    C^{(1^2)}P^{(1^2)}_{n=6} \ar[r,dash] & \dots
\end{tikzcd}
\]
\end{proposition}
\begin{proof}
    The $p=0,1$ vertices are by direct calculation. Also the first two arm vertices in each arm. And thereafter we may use 
\ref{eq:Cheby0} and in particular \ref{thm:main} and so on as explained above, 
noting that $\dim(\Deltaaa{n}{(n,\l)}) = 1$ here. 
\end{proof}

%\mdef 
%\ppm{[can we now DELETE this para?]}
%Now at the vertex $(p,\l)$ in the rollet graph (recall $p$ is the number of through lines, %\textit{i.e.} number of steps to the right, starting at 0 
%\ppm{at $\emptyset$}, 
%and $\l\vdash p_l$ is a partition labelling a Specht module for the appropriate symmetric group...), we place the gram determinant, $|\Delta_{(p,\l)}^{(p+4)}|$ (i.e. the two cup module labelled by the vertex), divided by the gram determinant of the modules appearing in its restriction, which includes $|\Delta_{(p-1,\l')}^{(p+3)}|$ (the two cup factors in the restriction) for all $\l'$ satisfying either $\l'=\l$ or which can be obtained from $\l$ by removing a box, and $|\Delta_{(p+1,\l'')}^{(p+3)}|$ (the one cup factors in the restriction) for all $\l'$ satisfying either $\l'=\l$ or which can be obtained from $\l$ by adding a box. Putting it mathematically: 

%Exactly as in (\ref{eq:MVF1}),  
%\ppmm{define the marginal vertex \sout{polynomials} functions by }
%\newcommand{\V}{{\mathcal V}}
%\[
%\V^{\ppm{n}}_{(p.\lambda)} \; = \; 
%\frac{ | \Delta_{(p,\l)}^{(n=p+2m)}| }{ \prod_{v\; adjacent\; to \; (p,\lambda)} |\Delta_{v}^{(n-1)}|}
%\]
%- then the next few figures are versions $\Roll^{(m)}_0$ of $\Roll_0$ where we put the marginal \bajm{[make the marginal a function of m not n so this is cleaner?]}
%for $n=p+2m$ for a fixed $m$ at each vertex $(p,\lambda)$.

\begin{proposition}\label{pr:l0}
(I) The graph $\Roll^{(1)}_0 $ 
(in which the $(p,\l)$ vertex is $\V^{p+2}_{(p,\l)}$)
is given by      %\ppm{[fix me!!--]}
\[ \hspace{-1.5cm} 
\begin{tikzcd}[row sep=0em,column sep=0.5em]
&&
\overbrace{
\tfrac{\alpha^{} \left(\alpha^2+\alpha -4\right)^{}}{(\alpha -1)^{} (\alpha +2)^{}} }^{ \frac{P^{(2)}_{4}}{P^{(2)}_{3}} }  \ar[r,dash]
  &
  \overbrace{
  \tfrac{\left(\alpha^4+\alpha^3-5 \alpha^2-\alpha +2\right)^{}}{\alpha^{} \left(\alpha ^2+\alpha -4\right)^{}}}^{ \V^{5}_{(3,(2))} = \frac{P^{(2)}_{5}}{P^{(2)}_{4}}} \ar[r,dash]
  &
  \overbrace{
  \tfrac{(\alpha -1)^{} \alpha^{} \left(\alpha ^3+2 \alpha ^2-4 \alpha -6\right)^{}}{\left(\alpha ^4+\alpha ^3-5 \alpha ^2-\alpha +2\right)^{}} }^{\frac{P^{(2)}_{6}}{P^{(2)}_{5}} \ppm{}}   \ar[r,dash]&
  \dots
\\
\alpha^{}\ar[r,dash] &
  \overbrace{
 (\alpha-1)^{} \tfrac{(\alpha-1)^{} (\alpha+2)^{}}{\alpha^{}}    }^{ \propto \frac{P^{(2)}_{3}}{P^{(2)}_{2}} \frac{P^{(1^2)}_{3}}{P^{(1^2)}_{2}}  }    \ar[ur,dash]\ar[dr,dash]
    \\
    &&
    \underbrace{
    \tfrac{(\alpha -2)^{} (\alpha +1)^{}}{(\alpha -1)^{}}  }_{ \frac{P^{(1^2)}_{4}}{P^{(1^2)}_{3}}  }  \ar[r,dash]
      &
      \underbrace{
    \tfrac{\left(\alpha ^3-\alpha ^2-3 \alpha +1\right)^{}}{(\alpha -2)^{} (\alpha +1)^{}} }_{\frac{P^{(1^2)}_{5}}{P^{(1^2)}_{4}}} \ar[r,dash]
    &
    \underbrace{
    \tfrac{(\alpha -1)^{} \left(\alpha ^3-4 \alpha -2\right)^{}}{\left(\alpha ^3-\alpha ^2-3 \alpha +1\right)^{}} }_{\frac{P^{(1^2)}_{\ppmm{6}}}{P^{(1^2)}_{5}}} \ar[r,dash]&\dots
\end{tikzcd}
\]
%that is         \[        \ppm{[same-but-using-Chebys-so-elipsis-is-clear]}        \]
%
(II) The graph $\Roll^{(2)}_0 $ \ppmm{(i.e. for $m=2$,  
so $n=4$ for vertex $\emptyset$; 
$n=5$ for vertex $(1)$; 
$n=6$ for $(2,(2))$ and $(2,(1^2))$; 
$n=7$ for $(3,(2))$ and $(3,(1^2))$; 
and so on)} 
starts (up to $p=8$):  % the graph $\Roll^{(2)}_{0}$ below:
\[  \hspace{-1.5cm} 
\begin{tikzcd}[row sep=0em,column sep=0.5em]
&&\frac{\alpha ^5 \left(\alpha ^2+\alpha -4\right)^5}{(\alpha -1)^5 (\alpha +2)^5}\ar[r,dash]
  &
  \overbrace{
\tfrac{\left(\alpha ^4+\alpha ^3-5 \alpha ^2-\alpha +2\right)^6}{\alpha ^6 \left(\alpha ^2+\alpha -4\right)^6} 
 }^{\V^7_{(3,(2))} = \left( \frac{ P^{(2)}_{5}}{P^{(2)}_{4}} \right)^{d^6_{(4,(2))}} }  \ar[r,dash]
  &
  \frac{(\alpha -1)^7 \alpha ^7 \left(\alpha ^3+2 \alpha ^2-4 \alpha -6\right)^7}{\left(\alpha ^4+\alpha ^3-5 \alpha ^2-\alpha +2\right)^7}\ar[r,dash]
  & \dots
  \\
    \alpha^3\ar[r,dash]&\frac{(\alpha -1)^8 (\alpha +2)^4}{\alpha ^4}\ar[ur,dash]\ar[dr,dash]
    \\
    &&
    \underbrace{
    \tfrac{(\alpha -2)^5 (\alpha +1)^5}{(\alpha -1)^5}   }_{ \left( \frac{ P^{(1^2)}_{4}}{P^{(1^2)}_{3}} \right)^{d^5_{(3,(1^2))}} }   \ar[r,dash]
      &
      \frac{\left(\alpha ^3-\alpha ^2-3 \alpha +1\right)^6}{(\alpha -2)^6 (\alpha +1)^6}\ar[r,dash]
      &
      \frac{(\alpha -1)^7 \left(\alpha ^3-4 \alpha -2\right)^7}{\left(\alpha ^3-\alpha ^2-3 \alpha +1\right)^7}\ar[r,dash]
      & \dots
\end{tikzcd}
\]
(III)
%\mdef 
%Now we do the same thing, but one step deeper, \textit{i.e.} three cup modules, divided by their restriction factors. That is, 
The graph $\Roll^{(3)}_0$ starts (up to $p=5$):
\[\begin{tikzcd}[row sep=0em,column sep=0.5em]&&\frac{\alpha ^{22} \left(\alpha ^2+\alpha -4\right)^{22}}{(\alpha -1)^{22} (\alpha +2)^{22}}\ar[r,dash]&\frac{\left(\alpha ^4+\alpha ^3-5 \alpha ^2-\alpha +2\right)^{29}}{\alpha ^{29} \left(\alpha ^2+\alpha -4\right)^{29}}\ar[r,dash]&\frac{(\alpha -1)^{37} \alpha ^{37} \left(\alpha ^3+2 \alpha ^2-4 \alpha -6\right)^{37}}{\left(\alpha ^4+\alpha ^3-5 \alpha ^2-\alpha +2\right)^{37}}\ar[r,dash]&\dots\\
    \alpha^{11}\ar[r,dash]&\frac{(\alpha -1)^{32} (\alpha +2)^{16}}{\alpha ^{16}}\ar[ur,dash]\ar[dr,dash]\\
    &&\frac{(\alpha -2)^{22} (\alpha +1)^{22}}{(\alpha -1)^{22}}\ar[r,dash]&\frac{\left(\alpha ^3-\alpha ^2-3 \alpha +1\right)^{29}}{(\alpha -2)^{29} (\alpha +1)^{29}}\ar[r,dash]&\frac{(\alpha -1)^{37} \left(\alpha ^3-4 \alpha -2\right)^{37}}{\left(\alpha ^3-\alpha ^2-3 \alpha +1\right)^{37}}\ar[r,dash]&\dots
\end{tikzcd}\]
\end{proposition}
\noindent 
{\em Proof}. 
(I) follows from \ref{pr:1cupgram}. 
\ppmm{(II,III) So far the Prop. is to be understood only to apply for the cases explicitly shown. Thus proof is by direct calculation.  \qed}

\ignore{\mdef
\bajm{[delete from here?]} 
\ppm{
Observe from Fig.\ref{fig:l=0Brat}, 
and cf. (\ref{eq:l=0+}), 
that the exponents follow (essentially) the exact same rule as in the $l=-1$ case, in (\ref{th:not-main--1}). 
The fine tuning is as follows. Firstly note that vertex $(1)$ is, in a sense, a new kind of case - since trivalent - and has non-trivial `edge factors' from two edges ... 
}

From this we have a conjecture for all $n$. 

\mdef 
The next steps (which can be done in either order!) will be:
\\
$\bullet\;$ Prove it.
\\
$\bullet\;$ Interrogate for the structure of standard modules - see \S\ref{ss:repthy}. 
\\
Then:
\\
$\bullet\;$  Lift to $l=1$ etc.!

\pagebreak

\mdef \bajm{[keep or not?]}
%\begin{landscape}
Let us have a look at $n=7$ row of determinants:
\[ \hspace{-2cm} 
\begin{tikzcd}
&
\baa{
(\alpha -2) (\alpha -1)^{10} \\ (\alpha +1) (\alpha +2) \\\left(\alpha ^3+2 \alpha ^2-4 \alpha -6\right) \\\left(\alpha ^4+\alpha ^3-5 \alpha ^2-\alpha +2\right)^7  \hspace{-1cm} \\ \ppm{(3,(2))}\\\bajm{22}
}    \hspace{-1cm} 
&
\baa{(\alpha -1) \\
\left(\alpha ^6+\alpha ^5-7 \alpha ^4-3 \alpha ^3+11 \alpha ^2+\alpha -2\right)\\\ppm{(5,(2))}\\\bajm{7}
}
&
\baa{1\\\ppm{(7,(2))}\\\bajm{1}
}
\\
\baa{
(\alpha -2)^8 (\alpha -1)^{60} (\alpha +1)^8\\ (\alpha +2)^{30} \left(\alpha ^2+\alpha -4\right)^8 \\\left(\alpha ^3-\alpha ^2-3 \alpha +1\right) \\\left(\alpha ^4+\alpha ^3-5 \alpha ^2-\alpha +2\right)\\\ppm{(1,(1))}\\\bajm{43}
}  \hspace{-1cm} 
\\
&
\baa{(\alpha -1)^{10} (\alpha +2)^8 \\ \left(\alpha ^2+\alpha -4\right)\\ \left(\alpha ^3-4 \alpha -2\right) \\ \left(\alpha ^3-\alpha ^2-3 \alpha +1\right)^7 \\ \ppm{(3,(1^2))}\\\bajm{22}
}  \hspace{-1cm} 
&
\baa{(\alpha -1) (\alpha +1) (\alpha +2)\\
\left(\alpha ^4-2 \alpha ^3-3 \alpha ^2+6 \alpha -1\right)\\\ppm{(5,(1^2))}\\\bajm{7}
}
&
\baa{1\\\ppm{(7,(1^2))}\\\bajm{1}
}
\end{tikzcd}\]
How to account for the factor $\left(\alpha ^2+\alpha -4\right)^8$ in $\Delttt{7}{1,(1)}$? 
The answer is now in Fig.\ref{fig:standardmapshenri}. 
%\end{landscape}

\medskip }

\subsubsection{Decorated Rollet Graph, for $l=1$} \label{ss:DRG1}
%Recall that the vertices of the Rollet graphs are pairs $(p,\lambda)$ where $p=0,1,\dots$ is the number of ``propagating lines" of the corresponding 
%\ppm{fibre of} modules, 
%and $\lambda $ is a partition of $\min(l+2,p)$. We organise the vertices into columns corresponding to the $p$ value (increasing from left to right) so, that we need only include the partition $\l$ as vertex label. 
%So for $l=1$ this is illustrated below.
Here we provide some examples of decorated Rollet graphs for $l=1$ to verify Theorem \ref{thm:V=C}. Firstly, recall the undecorated Rollet graph for $l=1$ is as follows:
\[
\begin{tikzcd}[column sep=0.6em,row sep=0em]
&&&(3)\ar[dash,r]&(3)\ar[dash,r]&(3)\ar[dash,r]&{}\\
&&(2)\arrow[dash,ur]\arrow[dash,dr]\\
\emptyset\arrow[dash,r] & (1)\arrow[dash,ur]\arrow[dash,dr] &&(2,1)\ar[dash,r]&(2,1)\ar[dash,r]&(2,1)\ar[dash,r]&{}\\
&&(1^2)\arrow[dash,ur]\arrow[dash,dr]\\
&&&(1^3)\ar[dash,r]&(1^3)\ar[dash,r]&(1^3)\ar[dash,r]&{}
\end{tikzcd}
\]
\ignore{Each node, then, labels an infinite family of $J_{1,n}$-modules for $n=p,p+2,p+4,\dots$. 
In the case of $n=p$, the corresponding module can simply be regarded as the 
\ppmm{image of the}
Specht module $\mathcal{S}^{\lambda}$, with any non-fully propagating diagram from $J_{l,n}$ acting as 0. 
For $n=p+2k$, we sometimes refer to the corresponding module as a $k$-cup module since a set of basis elements for the $J_{l,n}$-module can be described by pairs consisting of basis element for $\mathcal{S}^{\lambda}$, together with a half-diagram of height $\leq l$ with $k$-cups (\textit{cf.} \cite[Theorem something]{KadarYu})
\ppm{[-fix me!]}. 
We include the case $k=0$ in this terminology.

Now let us decorate this Rollet graph with the gram determinants of the $k$-cup modules (for $k>0$, since the $k=0$ case is not interesting \ppm{[LOL! you mean its a simple rep?]} over $\C$). 

\mdef 
\ppm{[explain a bit more!? this chunk seems to have been pulled from nowhere? linear combinations just correspond to tweaking the initial conditions, right?... Can this, then, be dragged to (or explained in terms of) the 2nd Cheby section, \S\ref{s:roots}?]}
To give some of the determinants here 
%For this, 
we use the following notation:
\begin{align*}
    F^{(3)}_{n}(\a)&=U_{n-2}(\a/2)+4U_{n-3}(\a/2)-U_{n-4}(\a/2)-2U_{n-5}(\a/2)-2U_{n-6}(\a/2)\\
    F^{(2,1)}_{n}(\a)&=U_{n-1}(\a/2)-4U_{n-3}(\a/2)-4U_{n-5}(\a/2)-2U_{n-7}(\a/2)\\
    F^{(1^3)}_{n}(\a)&=U_{n-3}(\a/2)-2U_{n-4}(\a/2)-2U_{n-5}(\a/2)
\end{align*}
\ppm{where $U_n(x)$ is as defined in \ref{where??} above.}

These polynomials (seem to  \ppm{[-need to deal with this somehow. ideally facts kept well away from speculations - in different sections.]}) satisfy the following:
\begin{itemize}
    \item $F^{(3)}_{n}(\a)$ is monic of degree $n-2$. It has $n-3$ roots of the form $z_k^{(n)}=2\cos (\th_k)$ for $\frac{(k-1)\pi }{n-3}<\th_k<\frac{k\pi }{n-3}$, for $k=1,\dots, n-3$, and one root of the form $z_0^{(n)}=-2 \cosh(\th_0)$ for $\th_0\in \R_{>0}$. Furthermore, $\lim_{n\to \infty} z_0^{(n)}$ seems to exist
    $\lim_{n\to \infty} z_0^{(n)}\approx -4.39356$.
    \item $F^{(1^3)}_{n}(\a)$ is monic of degree $n-3$. It has $n-4$ roots of the form $y_k^{(n)}=2\cos (\phi_k)$ for $\frac{k\pi }{n-3}<\phi_k<\frac{(k+1)\pi }{n-3}$, for $k=1,\dots, n-4$, and one root of the form $y_0^{(n)}=2 \cosh(\phi_0)$ for $\phi_0\in \R_{>0}$. Furthermore, $\lim_{n\to \infty} y_0^{(n)}$ seems to exist and $\lim_{n\to \infty} y_0^{(n)}\approx 3.09808$.
    \item $F^{(2,1)}_{n}(\a)$ is monic of degree $n-1$. It has $n-3$ roots of the form $x_k^{(n)}=2\cos (\psi_k)$ for $\frac{(2k-1)\pi }{2(n-2)}<\psi_k<\frac{(2k+1)\pi }{2(n-2)}$, for $k=1,\dots, n-3$, and two roots of the form $x_{0,\pm}^{(n)}=\pm 2 \cosh(\psi_0)$ for $\psi_0\in \R_{>0}$. Furthermore, $\lim_{n\to \infty} x_{0,+}^{(n)}$ seems to exist and $\lim_{n\to \infty} x_{0,+}^{(n)}\approx 2.66529$.
\end{itemize}
With this in mind, we proceed. 

\mdef 
The Rollet graph \ppm{[which one? gram? yes]} is given on the next page. 
We fill this graph with as much as we know for now - the $1$-cup gram det and the 2-cup one above this:
\ppm{[darn it, I hoped we'd dealt with the landscape problem. not sure yet what's achievable here. LOL, I guess this one is just even wider than before.]}}

\begin{landscape}
\begin{figure}
    \label{fig:placeholder}\[
\begin{tikzcd}[column sep=0.6em,row sep=0em]
&&&\begin{smallmatrix} \a (\a-2)^{8} (\a+2)^4(\a+4)\\
\times (\a^3+4\a^2-4\a-10){(\a^4-7\a^2+3)^2}\\
\times \left[(\a-2)^2\a^2(\a+1)(\a^2+3 \a-6)\right]^9
\\[0.5em] \hline  %[0.5em]
    (\a-2)^2\a^2(\a+1)(\a^2+3 \a-6)
\end{smallmatrix}\ar[dash,r]&\dots\ar[dash,r]&
\begin{smallmatrix}(\a+4)(\a+2)^4{(\a^4-7\a^2+3)^2}\\\times (\a-2)^8 \left[P^{(3)}_{p+3}(\a)\right]\\  \times\left[\a^2(\a-2)^2 \left[P^{(3)}_{p+2}(\a)\right]\right]^{p+6} \\[0.5em] \hline 
    \a^2(\a-2)^2 \left[P^{(3)}_{p+2}(\a)\right]
\end{smallmatrix}\ar[dash,r]&{}\\
&&\begin{smallmatrix}
    (\a-2)^{8} (\a+1)(\a+4) (\a+2)^{4}\\
    \times(\a^2+3\a-6){(\a^4-7\a^2+3)^2}\\
    \times \left[\a^3 (\a-2)^2 (\a+4)\right]^8\\[0.5em] \hline 
    \a^3(\a-2)^2(\a+4)
\end{smallmatrix}\arrow[dash,ur]\arrow[dash,dr]\\
\begin{smallmatrix}\a^3 \,(\a-1)^2  (\a+2)\\[0.5em] \hline 
\a\end{smallmatrix} \arrow[dash,r] & \begin{smallmatrix}(\a+4)(\a-2)^5(\a+2)^3\\\times \left[ (\a-1)^{2} (\a+2) \right]^7  \\[0.5em] \hline 
    (\a-1)^2(\a+2)
\end{smallmatrix}\arrow[dash,ur]\arrow[dash,dr] &&\begin{smallmatrix}(\alpha -3) (\alpha -2)^{42} (\alpha -1)^2  \alpha ^{12}\\ 
\times(\alpha +1)^4 (\alpha +2)^{17}(\alpha +4)^{12} \left(\alpha ^2-7\right)^2\\
\times\left(\alpha ^2+3 \alpha -6\right) \left(\alpha ^4-7 \alpha ^2+3\right)^{20}\\[0.5em] \hline
    (\alpha -2)^3 \alpha  (\alpha +2) (\alpha +4) {\left(\alpha ^4-7 \alpha ^2+3\right)^2}
\end{smallmatrix}\ar[dash,r]&\dots\ar[dash,r]&\begin{smallmatrix}
\frac{1}{\a^6}
    \left[\Delttt{(p+2)}{p,(2,1)}(\a)\right]^{p+6}
    \Delttt{(p+3)}{p+1,(2,1)}(\a)\Delttt{(4)}{2,(1^2)}(\a)\\
    \times \Delttt{(4)}{2,(2)}(\a)\Delttt{(5)}{3,(1^3))}(\a)\Delttt{(5)}{3,(3)}(\a)\Delttt{(5)}{3,(2,1)}(\a)\\[0.5em] \hline
    (\a-2)^3\a (\a+2)(\a+4) \left[P^{(2,1)}_{p+2}(\a)\right]^2
\end{smallmatrix}\ar[dash,r]&{}\\
&&\begin{smallmatrix}(\a-2)^{7}(\a+2)^{4}(\a+4)^2(\a+1)\\
\times\a (\a-3){(\a^4-7\a^2+3)^2}\\
\times \left[(\a-2)^3(\a+2)^3\right]^8\\[0.5em] \hline
    (\a-2)^3(\a+2)^3
\end{smallmatrix}\arrow[dash,ur]\arrow[dash,dr]\\
&&&\begin{smallmatrix}\a (\a-2)^{7} (\a+2)^{4}  (\a+4)^2 \\
\times(\a^3-2\a^2-4\a+2){(\a^4-7\a^2+3)^2}\\
\times \left[(\a-2)^2(\a-3 )(\a+1)(\a+2)^3\right]^9 \\[0.5em] \hline
    (\a-2)^2(\a-3 )(\a+1)(\a+2)^3
\end{smallmatrix}\ar[dash,r]&\dots\ar[dash,r]&
    \begin{smallmatrix} \a (\a+4)^2  {(\a^4-7\a^2+3)^2}\\
    \times (\a-2)^7(\a+2)^4 \left[P^{(1^3)}_{p+3}(\a)\right] \\
    \times \left[(\a-2)^2(\a+2 )^3 \left[P^{(1^3)}_{p+2}(\a)\right]\right]^{p+6}\\[0.5em] \hline
        (\a-2)^2(\a+2 )^3\left[ P^{(1^3)}_{p+2}(\a)\right]
    \end{smallmatrix}\ar[dash,r]&{}
\end{tikzcd}\]    
    \caption{Gram decorated rollet graph for $l=1$. We place two non-trivial determinants on each vertex (i.e. the cases $n=p+2,p+4$ on vertex $(p,\l)$). This general form of two cup determinants in the $\l$-arms for the row and column partitions is verified up to $p=7$.}
\end{figure}
\end{landscape}


%\ppm{[Caption this figure?:]}
\ignore{\bajm{[delete from here ***]}
Of course we know from the structure theorem for the Brauer algebra \cite{Martin09} that all the roots in the head part of the graph are integral  in these low ranks 
(indeed we know the gram determinants in the same way) until $n$ is large enough to have standard modules in the \arm s of the graph. 
Thus in the 1-cup Rollet section all the roots are integral until we reach the \hip: 
\ppm{[Caption this figure?:]}
\[
\begin{tikzcd}[column sep=0.6em,row sep=0em]
&&&\begin{smallmatrix} 
    (2)^2(0)^2(-1)\bajm{\dots}
\end{smallmatrix}\ar[dash,r]&\dots\ar[dash,r]&\begin{smallmatrix}
    (0)^2(2)^2 \left[F^{(3)}_{p+2}(\a)\right]
\end{smallmatrix}\ar[dash,r]&{}\\
&&\begin{smallmatrix}
    (0)^3(2)^2(-4)
\end{smallmatrix}\arrow[dash,ur]\arrow[dash,dr]\\
\begin{smallmatrix}(0)^1\end{smallmatrix} \arrow[dash,r] & \begin{smallmatrix}    (1)^2(-2)
\end{smallmatrix}\arrow[dash,ur]\arrow[dash,dr] &&\begin{smallmatrix}
    (2)^3 (0)  (-2) (-4) \bajm{\dots}
\end{smallmatrix}\ar[dash,r]&\dots\ar[dash,r]&\begin{smallmatrix}
    (-2)^3 (0) (-2)(-4) \left[F^{(2,1)}_{p+2}(\a)\right]^2
\end{smallmatrix}\ar[dash,r]&{}\\
&&\begin{smallmatrix}    (2)^3(-2)^3
\end{smallmatrix}\arrow[dash,ur]\arrow[dash,dr]\\
&&&\begin{smallmatrix}    (2)^2(3 )(-1)(-2)^3
\end{smallmatrix}\ar[dash,r]&\dots\ar[dash,r]&
    \begin{smallmatrix}         (2)^2(-2 )^3\left[ F^{(1^3)}_{p+2}(\a)\right]
    \end{smallmatrix}\ar[dash,r]&{}
\end{tikzcd}\] 

Let $\left|\Deltaaa{n}{(p,\l)}\right|$ denote the gram-det of the $J_{1,n}$-module labelled by the pair $(p,\l)$. We seem to have \ppm{[-time to adjust this phrasing.]} observed the following factorisation rules for $2$-cup determinants in terms of $1$-cup determinants:
\begin{align*}
 \left|\Delta_{(p,(1^3))}^{(p+4)}\right|&=\frac{1}{\a^3}
    \left|\Delta_{(p,(1^3))}^{(p+2)}\right|^{p+6}
    \left|\Delta_{(p+1,(1^3))}^{(p+3)}\right|
    \left|\Delta_{(2,(2))}^{(4)}\right|
    \left|\Delta_{(3,(2,1))}^{(5)}\right|\\
 \left|\Delta_{(p,(3))}^{(p+4)}\right|&=\frac{1}{\a^3}
    \left|\Delta_{(p,(3))}^{(p+2)}\right|^{p+6}
    \left|\Delta_{(p+1,(3))}^{(p+3)}\right|
    \left|\Delta_{(2,(1^2))}^{(4)}\right|
    \left|\Delta_{(3,(2,1))}^{(5)}\right|\\
 \left|\Delta_{(p,(2,1))}^{(p+4)}\right|&=\frac{1}{\a^6}
    \left|\Delta_{(p,(2,1))}^{(p+2)}\right|^{p+6}
    \left|\Delta_{(p+1,(2,1))}^{(p+3)}\right|
    \left|\Delta_{(2,(1^2))}^{(4)}\right|\left|\Delta_{(2,(2))}^{(4)}\right|\\
    &\quad \times\left|\Delta_{(3,(1^3))}^{(5)}\right|\left|\Delta_{(3,(3))}^{(5)}\right|   
    \left|\Delta_{(3,(2,1))}^{(5)}\right|\\
\end{align*}

\bajm{[to here ***]}}

\mdef \label{pa:l1}
We now include 
(by direct computation)
the \ppmm{section $\Roll^{(2)}_{1}$ through the }
decorated Rollet graph $\RollV_1$ as per the $l=0$ case in \ref{pr:l0} II. This computation completes the cases remaining in Theorem \ref{thm:V=C} (the $p=7$ case is omitted since it did not fit in the margins).
%\rotatebox{45}{
%\begin{landscape}
\[
\hspace{-4.51cm} 
\begin{tikzcd}[column sep=0.6em,row sep=0em]
&&&\frac{(\alpha +1)^8 \left(\alpha ^2+3 \alpha -6\right)^8}{\alpha ^8 (\alpha +4)^8}\ar[dash,r]
   &
   \overbrace{
   \tfrac{\alpha ^9 \left(\alpha ^3+4 \alpha ^2-4 \alpha -10\right)^9}{(\alpha +1)^9 \left(\alpha ^2+3 \alpha -6\right)^9}   }^{ \left( \frac{ P^{(3)}_{\ppmm{6}}}{P^{(3)}_{5}} \right)^9}   \ar[dash,r]
   &
   \overbrace{
   \tfrac{\left(\alpha ^5+4 \alpha ^4-5 \alpha ^3-14 \alpha ^2+3 \alpha +6\right)^{10}}{\alpha ^{10} \left(\alpha ^3+4 \alpha ^2-4 \alpha -10\right)^{10}}   }^{  \left( \frac{ P^{(3)}_{\ppmm{7}}}{P^{(3)}_{6}} \right)^{10}  }   \ar[dash,r]
   &\dots
   \\
&&\frac{(\alpha -2)^{14} \alpha ^{14}}{(\alpha -1)^{14}} \frac{\a^7(\alpha +4)^7}{(\alpha +2)^7}\arrow[dash,ur]\arrow[dash,dr]
\\
\alpha ^3\arrow[dash,r] & (\alpha -1)^{6}\frac{(\alpha -1)^{6} (\alpha +2)^6}{\alpha ^6}\arrow[dash,ur]\arrow[dash,dr] 
  &&
  \underbrace{
  \tfrac{\left(\alpha ^4-7 \alpha ^2+3\right)^{16}}{(\alpha -2)^{16} \alpha ^{16} (\alpha +2)^{16}} }_{  \left( \frac{ P^{(21)}_{\ppmm{5}}}{P^{(21)}_{4}} \right)^{16}     } 
    \ar[dash,r]
  &
  \frac{(\alpha -1)^{18} \alpha ^{18} (\alpha +1)^{18} \left(\alpha ^2-7\right)^{18}}{\left(\alpha ^4-7 \alpha ^2+3\right)^{18}}\ar[dash,r]&\frac{\left(\alpha ^6-9 \alpha ^4+14 \alpha ^2-3\right)^{20}}{(\alpha -1)^{20} \alpha ^{20} (\alpha +1)^{20} \left(\alpha ^2-7\right)^{20}}\ar[dash,r]&\dots
\\
&&(\alpha -2)^{7}\frac{(\alpha -2)^{14} (\alpha +2)^{14}}{(\alpha -1)^{14}}\arrow[dash,ur]\arrow[dash,dr]
\\
&&&
  \underbrace{
  \tfrac{(\alpha -3)^8 (\alpha +1)^8}{(\alpha -2)^8} }_{ \left( \frac{ P^{(1^3)}_{\ppmm{5}}}{P^{(1^3)}_{4}} \right)^{8}    }     \ar[dash,r]
  &
  \frac{\left(\alpha ^3-2 \alpha ^2-4 \alpha +2\right)^9}{(\alpha -3)^9 (\alpha +1)^9}\ar[dash,r]&\frac{\left(\alpha ^4-2 \alpha ^3-5 \alpha ^2+4 \alpha +3\right)^{10}}{\left(\alpha ^3-2 \alpha ^2-4 \alpha +2\right)^{10}}\ar[dash,r]&\dots
\end{tikzcd}
\]    
%\end{landscape}
%}
Where the $P_{n}^{(\l)}$ are as computed in \S\ref{ss:examples0} (in particular, \textit{cf.} \ref{pa:1dgdets} and \ref{pa:21dets}).

Observe that
\[
\CC^{6}_{(2,(2))} = 
\left( \frac{ P^{(21)}_{\ppmm{4}}}{P^{(21)}_{3}} \right)^{14} 
\left( \frac{ P^{(3)}_{\ppmm{4}}}{P^{(3)}_{3}} \right)^{7} 
=  \left( \frac{\xx(\xx^2 -4)}{3(\xx^2 -1)}  \right)^{14} 
\left( \frac{\xx (\xx+4)}{3(\xx+2)}  \right)^{7} 
\;\propto \; (\xx+2)^{14} \; \V^6_{(2,(2))}
\]
\[
\CC^{6}_{(2,(1^2))} =
\left( \frac{ P^{(21)}_{\ppmm{4}}}{P^{(21)}_{3}} \right)^{14} 
\left( \frac{ P^{(1^3)}_{\ppmm{4}}}{P^{(1^3)}_{3}} \right)^{7} 
=  \left( \frac{\xx(\xx^2 -4)}{3(\xx^2 -1)}  \right)^{14} 
\left( \frac{ (\xx-2)}{3}  \right)^{7} 
\;\propto  \; \xx^{14} \; \V^6_{(2,(1^2))}
\]
so the \ournotmain\ conjecture narrowly fails here to extend into the head.
%\\
%\bajm{[what?]}One can compare this with \ref{pa:Cheby1}. 



\subsubsection{Decorated Rollet Graph, for $l=2$}   \label{ss:gdetl2}

We now continue the same process for $l=2$.
The (undecorated) Rollet graph in integer-partition (as opposed to Young diagram) notation is:
\[
\begin{tikzcd}[column sep=0.6em,row sep=0em]
&&&&(4)\ar[dash,r]&(4)\ar[dash,r]&(4)\ar[dash,r]&\dots\\
&&&(3)\ar[dash,ur]\ar[dash,dr]\\
&&(2)\arrow[dash,ur]\arrow[dash,dr]&& (3,1)\ar[dash,r]&(3,1)\ar[dash,r]&(3,1)\ar[dash,r]&\dots\\
\emptyset\arrow[dash,r] & (1)\arrow[dash,ur]\arrow[dash,dr] &&(2,1)\ar[dash,ur]\ar[dash,dr]\ar[dash,r]&(2,2)\ar[dash,r]&(2,2)\ar[dash,r]&(2,2)\ar[dash,r]&\dots\\
&&(1^2)\arrow[dash,ur]\arrow[dash,dr]&&(2,1^2)\ar[dash,r]&(2,1^2)\ar[dash,r]&(2,1^2)\ar[dash,r]&\dots\\
&&&(1^3)\ar[dash,dr]\ar[dash,ur]\\
&&&&(1^4)\ar[dash,r]&(1^4)\ar[dash,r]&(1^4)\ar[dash,r]&\dots
\end{tikzcd}
\]
In Figure \ref{fig:l2} we give the section $\Roll_2^{(2)}$ in all cases so far determined by direct computation. 
\begin{figure}
\[
\hspace{-3.51cm} \begin{tikzcd}[column sep=0.6em,row sep=0em]
&&&&\frac{(\alpha +2)^{12} \left(\alpha ^2+5 \alpha -8\right)^{12}}{(\alpha +1)^{12} (\alpha +6)^{12}}\ar[dash,r]&\frac{\left(\alpha ^4+7 \alpha ^3+\alpha ^2-23 \alpha -6\right)^{13}}{(\alpha +2)^{13} \left(\alpha ^2+5 \alpha -8\right)^{13}}\ar[dash,r]&{}\\
&&&\frac{(\alpha -2)^{33} (\alpha +1)^{44} (\alpha +6)^{11}}{\alpha ^{33} (\alpha +4)^{11}}\ar[dash,ur]\ar[dash,dr]\\
&&\frac{(\alpha -2)^{20} \alpha ^{30} (\alpha +4)^{10}}{(\alpha -1)^{20} (\alpha +2)^{10}}\arrow[dash,ur]\arrow[dash,dr]&& (\dots)\ar[dash,r]&(\dots)\ar[dash,r]&{}\\
\alpha ^3 \arrow[dash,r] & \frac{(\alpha -1)^{12} (\alpha +2)^6}{\alpha ^6} \arrow[dash,ur]\arrow[dash,dr] &&\frac{ (\alpha -3)^{55} (\alpha -1)^{66} (\alpha +1)^{55} (\alpha +4)^{33}}{ (\alpha -2)^{55} \alpha ^{33} (\alpha +2)^{33}}\ar[dash,ur]\ar[dash,dr]\ar[dash,r]&(\dots)\ar[dash,r]&(\dots)\ar[dash,r]&{}\\
&&\frac{(\alpha -2)^{30} (\alpha +2)^{20}}{(\alpha -1)^{20}}\arrow[dash,ur]\arrow[dash,dr]&&(\dots)\ar[dash,r]&(\dots)\ar[dash,r]&{}\\
&&&\frac{(\alpha -3)^{44} (\alpha +2)^{33}}{(\alpha -2)^{33}}\ar[dash,dr]\ar[dash,ur]\\
&&&&\frac{(\alpha -4)^{12} (\alpha +1)^{12}}{(\alpha -3)^{12}}\ar[dash,r]&\frac{\left(\alpha ^3-3 \alpha ^2-5 \alpha +3\right)^{13}}{(\alpha -4)^{13} (\alpha +1)^{13}}\ar[dash,r]&{}
\end{tikzcd}
\]
\caption{The section, $\Roll_2^{(2)}$, of the decorated Rollet graph, $\RollV_2$, in all cases so far computed. Any entry marked $(\dots)$ is so far not determined by direct computation.
\label{fig:l2}}
\end{figure}
 
Observe that
\begin{align*}
\CC^{7}_{(3,(3))} = 
\left( \frac{ P^{(31)}_{\ppmm{5}}}{P^{(31)}_{4}} \right)^{33} 
\left( \frac{ P^{(4)}_{\ppmm{5}}}{P^{(4)}_{4}} \right)^{11} 
&=  \left( \frac{(\alpha -2) (\alpha -1) (\alpha +1) (\alpha +4)}{4 \alpha  \left(\alpha ^2+\alpha -4\right)}  \right)^{33} 
\left( \frac{(\alpha +1) (\alpha +6)}{4 (\alpha +4)}\right)^{11} 
\\
&\propto \; \frac{\left(\alpha ^2+\alpha -4\right)^{33}}{(\alpha -1)^{33} (\alpha +4)^{33}} \; \V^7_{(3,(3))}
\end{align*}
\begin{align*}
\CC^{7}_{(3,(2,1))} = 
\left( \frac{ P^{(31)}_{\ppmm{5}}}{P^{(31)}_{3}} \right)^{33} 
\left( \frac{ P^{(4)}_{\ppmm{5}}}{P^{(4)}_{4}} \right)^{11} 
&= \left( \frac{(\alpha -2) (\alpha -1) (\alpha +1) (\alpha +4)}{4 \alpha  \left(\alpha ^2+\alpha -4\right)}  \right)^{33}\\
&\quad \times \left(\frac{(\alpha -3) (\alpha -1) (\alpha +2)}{4 (\alpha -2) (\alpha +1)}\right)^{33}\left(\frac{(\alpha -3) (\alpha +1)}{4 (\alpha -2)}\right)^{22}\\ 
\\
&\propto \; \frac{ (\alpha +1)^{33} \left(\alpha ^2+\alpha -4\right)^{33}}{(\alpha -2)^{33} (\alpha +2)^{66}} \; \V^7_{(3,(2,1))}
\end{align*}
\begin{align*}
\CC^{7}_{(3,(1^3))} = 
\left( \frac{ P^{(31)^T}_{\ppmm{5}}}{P^{(31)^T}_{4}} \right)^{33} 
\left( \frac{ P^{(4)^T}_{\ppmm{5}}}{P^{(4)^T}_{4}} \right)^{11} 
&=  \left( \frac{(\alpha -3) (\alpha -1) (\alpha +2)}{4 (\alpha -2) (\alpha +1)}  \right)^{33} 
\left( \frac{\alpha -3}{4}\right)^{11} 
\\
&\propto \; \frac{(\a+1)^{33}}{(\a-1)^{33}} \; \V^7_{(3,(1^3))}
\end{align*}
which again suggests the \ournotmain\ conjecture can be modified suitable to extend into the head.